%% QDCVR platform — Knowledge-Based Systems submission (elsarticle)
%% v2 restructure 2026-10-12: content-retrieval-based KB *management system* framing.
%% Compile: pdflatex main2 && bibtex main2 && pdflatex main2 && pdflatex main2
%% AUTHOR ACTION: replace author placeholders, ORCID, and funding fields.
\documentclass[preprint,12pt]{elsarticle}

\usepackage[T1]{fontenc}
\usepackage[utf8]{inputenc}
\usepackage{booktabs}
\usepackage{array}
\usepackage{graphicx}
\usepackage{amsmath}
\usepackage{amssymb}
\usepackage{xcolor}
\usepackage{url}
\usepackage[colorlinks=true,linkcolor=blue,citecolor=blue,urlcolor=blue]{hyperref}
\usepackage{xurl}
\usepackage{caption}
\captionsetup{font=small,labelfont=bf}
\newcommand{\yes}{\checkmark}
\newcommand{\no}{--}

\graphicspath{{figs/}}
\setlength{\emergencystretch}{1.5em}

%% Float placement: the paper carries nine figures and twenty-one table
%% environments (five in the appendix: one per-question rank table and four
%% supplementary tables), so the
%% default fractions defer most of them past the references. Allow denser
%% float pages to keep each float near the text that discusses it.
\renewcommand{\topfraction}{0.92}
\renewcommand{\bottomfraction}{0.85}
\renewcommand{\textfraction}{0.06}
\renewcommand{\floatpagefraction}{0.70}
\setcounter{topnumber}{4}
\setcounter{bottomnumber}{3}
\setcounter{totalnumber}{5}
\setcounter{dbltopnumber}{3}

\journal{Knowledge-Based Systems}

\begin{document}

\begin{frontmatter}

\title{Design of a content-retrieval-based knowledge-base management system
for trustworthy enterprise document intelligence}

\author[inst1]{Author Name\corref{cor1}}
\ead{author@example.com}
\author[inst1]{Co-author Name}
\cortext[cor1]{Corresponding author.}
\affiliation[inst1]{organization={Affiliation Department, University,
City, Country}, city={City}, country={Country}}

\begin{abstract}
Enterprises accumulate regulations, failure analyses and operational lessons
as unstructured documents, and answering one engineering question requires
locating, reading and citing the right one. Retrieval-augmented generation
addresses the locating step, but its ranking signal measures topic relatedness
rather than evidence presence. QDCVR is a self-hosted
knowledge-base system governing content across all three
actions. Six subsystems---quality-gated parsing and ingestion, content-routed
organization, life-cycle management with three-source consistency audits, a
document graph, an operational-knowledge store with credibility tiers, and
multi-lane content-verified retrieval---expose 95 Model Context Protocol tools
over one storage layer, each terminating at an explicit contract. We evaluate it on a deployed industrial corpus, a
322-document organization campaign, a 100-paper acquisition campaign and a
cross-lingual corpus. Twenty-one of 100 submitted
papers never became retrievable; the accounting gap is reported, not hidden. The gate exposes a
same-domain score-inflation failure consistent with LLM-judge
miscalibration: non-gold scores are indistinguishable from gold
(Mann--Whitney $p = 0.26$), judged re-ranking lowers Hit@5 from 0.94 to 0.44,
while reusing the same verdicts as a filter restores 0.94 (sign test
$p = 0.008$). Component ablations return a largely null result: on nine of ten
questions both arms received identical evidence, so verification's measured
value here is fail-closed provenance rather than ranking, and the losses sit
in the reading layer. Functional campaigns confirm ingestion,
deduplication, organization and abstention end to end; experience distillation
is reported as case evidence. Engine
choice moves cost by one to two orders of magnitude. Verifier calibration is
the open problem, and honesty needs an enforcement layer.
\end{abstract}

\begin{keyword}
Knowledge base management \sep Document ingestion \sep Content verification \sep
Knowledge organization \sep Retrieval-augmented generation \sep
Model Context Protocol
\end{keyword}

\end{frontmatter}

%% ============================================================ body (v2 order)
\section{Introduction}
\label{sec:intro}

Enterprises accumulate process regulations, failure analyses, and operational
lessons as unstructured documents on shared drives. Answering a concrete
engineering question --- \emph{why did edge delamination appear in an
electroluminescence inspection of a laminated photovoltaic module?} ---
requires locating the one report among thousands that records the decision and
reading it. Retrieval-augmented generation (RAG) addresses the location step by
embedding documents and the query into a shared vector space
\cite{rag-lewis,dpr}, and agentic stacks let a large language model orchestrate
retrieval tools iteratively \cite{mcp,singh-agentic-rag}. The dominant ranking
signal, dense cosine similarity, measures \emph{topic relatedness} rather than
\emph{evidence presence} \cite{siren,gao-ragsurvey}.

Three actions this task requires are addressed by three research communities,
each solving its own and leaving the others' problems intact. Retrieval
research optimizes \emph{finding}: sparse, dense and hybrid rankers, and more
recently graph-based and hierarchical indexes, are judged by whether a relevant
passage appears in a top-$k$ list. Knowledge management addresses
\emph{organizing}: classification schemes, controlled vocabularies and
lessons-learned systems organize how documents are arranged and reused, but do
not close the loop to an answer a machine must justify. Document AI
addresses \emph{reading}: layout-aware parsing recovers text, tables and
formulas from scanned pages \cite{mineru}, but a parsed document still has to
be placed, indexed, governed and trusted. In deployment the three concerns
compound.

The object all three actions govern is \emph{content}: a knowledge-base
management system expresses that position as enforceable contracts rather than
as good practice. Three commitments follow, each implemented and measured here.
First, \emph{content decides placement}: library membership, tags, descriptions
and graph edges are derived from what a document says, not from the folder it
arrived in, and the resulting catalog is the first retrieval instrument.
Second, \emph{content decides admissibility}: no similarity score, however
high, promotes a candidate into an answer, because the platform must itself
fetch the candidate's content and judge it to contain direct evidence. Third,
\emph{absence of evidence is an answer}: when verification eliminates every
candidate, the system returns an explicit not-found report rather than the
closest text. Claims are therefore auditable, with every placement, verdict,
cut and lifecycle operation on record.

QDCVR (Query-Driven Content-Verified Retrieval, v2.3), the self-hosted
knowledge-base management system built here,
implements this position end to end on a deployed industrial corpus and on
acquisition and functional campaigns. Six functional subsystems carry the
platform: a quality-gated parsing and ingestion pipeline, knowledge
organization with content routing, tags and a document graph, life-cycle
management with three-source consistency audits, an operational-knowledge store
with credibility tiers, a multi-lane content-verified retrieval layer, and a
tool interface exposing the platform as 95 Model Context Protocol tools
\cite{mcp}; the system runs on a single workstation. On the acquisition side,
of 100 papers submitted to a 100-document campaign, 21 never became
retrievable, and the
design decision that matters is that a document which fails a gate is
reported rather than
stored as an empty document (Section~\ref{sec:e1}). On the retrieval side, the
content-verification gate is a better filter than a ranker: reusing the
\emph{same} verdicts as a re-ranking signal lowers Hit@5 from 0.94 to 0.44,
while using them as a binary filter restores 0.94 (sign test $p = 0.008$),
because same-domain score inflation makes non-gold candidates statistically
indistinguishable from gold ones (Mann--Whitney $p = 0.26$)
(Section~\ref{sec:eval-verify}).

\noindent Contributions:
\begin{enumerate}
\item \textbf{A content-governed knowledge-base management architecture}
(Sections~\ref{sec:system}--\ref{sec:implementation}): requirements derived
from deployment, six design principles enforced in code, a four-layer
architecture, and an explicit data model whose contracts --- parse gates,
placement rules, three-source consistency, promotion conditions, task
semantics --- are stated as obligations rather than conventions.
\item \textbf{A quality-gated acquisition pipeline with a measured failure
budget} (Section~\ref{sec:acquisition}, E1): multi-format parsing with OCR,
exact and near-duplicate detection, structure-aware splitting of oversized
documents with part addressing, content-derived descriptions and tags, and an
index post-condition; evaluated on a 100-document campaign (79 retrievable,
a 21\% retrievability gap whose cause the campaign record does not
attribute) and on a cross-lingual corpus.
\item \textbf{Content-routed organization and life-cycle management}
(Section~\ref{sec:organization}, E2): placement by content, catalog and
metadata registries with progressive disclosure, tag governance, duplicate
hygiene, a document graph with an explicitly restricted role, an experience
store with credibility tiers, and lifecycle operations certified by a
read-only audit (322 documents in five libraries).
\item \textbf{Multi-channel content-verified retrieval} (Section~\ref{sec:lanes},
E4--E5): three lanes over one tool surface under enforced budgets, a
verification gate served by a resident judging worker, and an honest
abstention path --- compared against nine retrieval configurations, two of
them derived offline from the same logged candidate pools, and against a
bare agent and a pure-dense controller. What the ablation measures is where
the verdict belongs, and what the platform delivers is an auditable record of
every candidate and every rejection, not a recall gain.
\item \textbf{A verifier-failure analysis and its role decomposition}
(Section~\ref{sec:eval-verify}, E4, E6): same-domain inflation quantified
inside a deployed pipeline, the filter/re-ranker asymmetry, a threshold that
never abstains out of corpus, and robustness of the failure to paraphrase, to
a second verification signal and to a second corpus.
\item \textbf{An operational account of the system} (Sections~\ref{sec:impl}
and~\ref{sec:discussion}, E7): resident-worker and subprocess engineering,
engine-agnostic harness integration with cost measured across three engines,
a catalogue of failure modes with their detection symptoms, and the
non-functional behaviour observed under concurrent writes.
\end{enumerate}

A demo describing the platform's organization and retrieval layers was
published as a CIKM~2026 demo \cite{qdcvr-demo}. This paper extends it with
four components the demo does not contain: the acquisition pipeline with its
measured retrievability gap and cross-lingual evaluation; life-cycle
management and integrity auditing as a designed subsystem rather than a
maintenance procedure; the operational-knowledge store with its credibility
tiers and distillation loop; and the verifier-failure analysis with the
filter/re-ranker decomposition. The demo also reports no negative result, which
is the centrepiece here.

\subsection{Research questions and reading guide}
\label{sec:rq}

The paper is organized as a design-to-evidence arc around five research
questions, answered by one evidence group each and reported in one place
(Table~\ref{tab:ataglance}).

\begin{description}
\item[RQ1 (Enforceable contracts).] How can the three content commitments be
turned into system contracts that hold across retrieval strategies and
independent agent harnesses, rather than into prompt conventions?
\item[RQ2 (Acquisition and organization at scale).] Do quality-gated
acquisition and content-routed organization produce collections that are
usable and internally consistent when they grow, and what does the gate
reject?
\item[RQ3 (Content-verified retrieval).] Does staged, multi-channel retrieval
with content verification outperform one-shot vector top-$k$ and unverified
fusion --- and by which mechanism?
\item[RQ4 (Cost and benefit of verification).] What does the gate cost and
buy, and is its failure mode structural rather than an artefact of one
checkpoint, domain or language?
\item[RQ5 (Operational contract).] Where do honest abstention, latency, cost
and stability actually live in a running system?
\end{description}

\begin{table}[!htbp]
\centering
\footnotesize
\setlength{\tabcolsep}{3pt}
\caption{This paper at a glance: each research question, the evidence group
that answers it, and the headline finding.}
\label{tab:ataglance}
\begin{tabular}{p{0.8cm}p{3.7cm}p{2.9cm}p{4.9cm}}
\toprule
RQ & Question & Evidence & Headline finding \\
\midrule
RQ1 & Are the content commitments enforceable contracts? &
Architecture, interface, implementation
(\S\ref{sec:system}--\S\ref{sec:implementation}) & Yes: gates, placement
rules, lane whitelists and fail-closed verdicts make them properties of the
enforcement layer \\
RQ2 & Do gated acquisition and content routing scale? &
E1, E2, E8 (\S\ref{sec:acquisition}, \S\ref{sec:organization}) & Partly: routing
and audits hold at 322 documents, while 21 of 100 submitted papers never
became retrievable and the record does not attribute the gap. Routing is
measured as coverage rather than accuracy, because every stored filename
carries the acquisition prefix (\S\ref{sec:e8}) \\
RQ3 & Does staged verified retrieval beat one-shot top-$k$? &
E4 (\S\ref{sec:eval-retrieval}) & Yes on reachability, through staged
recall plus content adjudication; two-stage retrieval is an accelerator, not
the contribution \\
RQ4 & What does the gate cost and buy, and is the failure structural? &
E4, E6 (\S\ref{sec:eval-verify}) & Costs $<$1\,s; as a re-ranker it loses
Hit@5 ($0.94\!\to\!0.44$), as a filter it is neutral ($0.94$,
$p=0.008$); the loss tracks confuser density, not checkpoint or domain \\
RQ5 & Where do honesty, latency and cost live? &
E5, E7 (\S\ref{sec:eval-agent}, \S\ref{sec:implementation}) & Honesty lives
in the reading and answering layer, never in the threshold; latency is
dominated by LLM turns, and engine choice moves cost by two orders of
magnitude \\
\bottomrule
\end{tabular}
\end{table}

The remainder of the paper follows the arc.
Section~\ref{sec:related} positions the work against knowledge management,
document parsing, retrieval, verification and agent protocols.
Sections~\ref{sec:system}--\ref{sec:organization} present the architecture,
the acquisition pipeline and the organization subsystem, answering RQ1 and
RQ2. Section~\ref{sec:lanes} presents the retrieval lanes and the verification
gate, walking one real question end to end, and states what the retrieval layer
claims (RQ3). Section~\ref{sec:implementation} reports the engineering that
makes the contracts enforceable. Section~\ref{sec:eval} carries the eight
evidence groups (RQ4--RQ5).
Sections~\ref{sec:discussion}--\ref{sec:conclusion} distil the lessons and
state the threats to validity. Table~\ref{tab:modules} maps system functions
to components and evidence; Table~\ref{tab:ataglance} maps research questions
to evidence.

\section{Related work}
\label{sec:related}

\subsection{Knowledge bases, knowledge organization, and knowledge management}
Large public knowledge bases such as DBpedia and YAGO extract structured facts
from semi-structured sources under an ontology
\citep{auer-dbpedia-kb-2007,hoffart-yago2-kb-2013}. The enterprise content
management literature treats organizational content as a governance problem
rather than an extraction problem \citep{paivarinta-ecm-2005}; later work
characterizes the field's research agenda \citep{tyrvaainen-ecm-2006}.
Organizational knowledge creation describes how tacit experience becomes
explicit and shareable through socialization, externalization, combination, and
internalization, sustained by a shared context
\citep{nonaka-seci-1994,nonaka-seci-ba-2000}. Complementary strands study
metadata quality in digital repositories \citep{park-metadata-quality-2009}, the
emergent vocabulary of collaborative tagging \citep{golder-tagging-2006},
human-in-the-loop review of machine-produced labels
\citep{mosqueira-rey-hitl-2022}, and the auditability of automated decisions
\citep{barredo-arrieta-xai-2020}.

That literature names inconsistent metadata, uncontrolled vocabulary, unverified
labels, and unreusable experience, but treats them as management and social
problems and rarely closes the loop to machine answering. Closing that loop is
the contribution here: the content-derived metadata that organizes a base also
constrains retrieval, verification, and abstention
(Sections~\ref{sec:organization}--\ref{sec:gate}).

\subsection{Document parsing and acquisition pipelines}
Document AI supplies the components of an acquisition pipeline. PubLayNet
established large-scale layout analysis \citep{zhong-publaynet-layout-2019},
unified text--image pre-training improved layout-aware understanding of tables
and figures \citep{huang-layoutlmv3-docai-2022}, and MinerU provides the open
extraction pipeline for scientific and scanned documents that this platform
drives \citep{mineru}. Near-duplicate detection by locality-sensitive hashing
over shingles is the redundancy-control mechanism
\citep{manku-near-dup-2007}, and provenance has mature foundations in the W3C
PROV family \citep{missier-prov-2013} and in database provenance
\citep{cheney-provenance-2009}.

Parsing is usually presented as preprocessing, yet an extraction error
propagates into chunking, description, indexing, and every later citation and is
invisible in retrieval metrics. In our campaign 21 of 100 submitted arXiv papers
never reached a retrievable state (Section~\ref{sec:acquisition}), so
acquisition itself carries a quality gate, duplicate check, index-verification
step, and per-document provenance record.

\subsection{Retrieval: sparse, dense, hybrid, and RAG}
Retrieval-augmented generation grounds generation in retrieved passages
\citep{rag-lewis}; dense retrievers learn a shared query--passage space
\citep{dpr}; multi-granularity embedders such as BGE-M3 combine dense and sparse
signals \citep{bge-m3}. Lexical probabilistic ranking remains a strong baseline
\citep{bm25}, reciprocal-rank fusion merges ranked lists without training
\citep{rrf-cormack}, and ranking is benchmarked on MS~MARCO
\citep{craswell-msmarco-2021}. Graded-relevance metrics are the evaluation
instrument \citep{jarvelin-cumulated-gain-2002}, and surveys systematize the
design space as retrieve-then-read pipelines \citep{gao-ragsurvey}. Vectors live
in a persistent store with graph-based approximate nearest-neighbour search
\citep{chroma,hnsw}.

The lexical baseline is not the weak part: on the gold-annotated corpus BM25
reaches Hit@5 of 1.00. This paper does not compete on ranking quality. It treats
every such signal as a candidate generator that must still be confirmed against
content (Section~\ref{sec:lanes}).

\subsection{Structured and graph-based organization}
GraphRAG builds entity--relation graphs and community summaries for
query-focused summarization \citep{graphrag}; LightRAG reduces the idea to
dual-level retrieval \citep{lightrag}. RAPTOR recursively abstracts document
trees \citep{raptor}; the acquisition, representation, and application taxonomy
\citep{ji-kg-survey-2022} joins knowledge-graph work on embedding-based
representation and link prediction \citep{wang-kg-embedding-2017}, with
property-graph stores as the substrate \citep{neo4j}. Structured organization
also mitigates ungrounded generation \citep{siren}.

Our graph layer is deliberately thinner: Document, KnowledgeBase, and Tag nodes
joined by membership, tag, and weighted relatedness edges, built from metadata
that acquisition already verified. It navigates and expands candidates and never
answers alone, because keyword-limited graph search is the weakest standalone
signal on our corpus (Hit@5 0.56, rank quality MRR 0.562; the same run
appears as Table~\ref{tab:retrieval}'s graph row). It is therefore exposed as one channel among
several (Section~\ref{sec:lanes}).

\subsection{Content verification, correction, and abstention}
Self-RAG trains a generator to critique retrieved evidence and its own output
\citep{selfrag}; CRAG adds a corrective retriever that re-evaluates and
re-queries when retrieval quality is low \citep{crag}; FLARE interleaves
generation with forward-looking retrieval triggers \citep{flare};
Kamath et al.~\citep{kamath} formalize selective prediction under domain shift; CyberBOT
applies ontology-grounded retrieval in a reliability-critical domain
\citep{cyberbot}. Verifiability measurement in search
\citep{liu-verifiability} and citation-capable generation \citep{gao-citations}
frame the problem this paper attacks: citations that are present but not
supportive.

\emph{LLM-as-a-judge.} The verification gate is a locally hosted judgement model
\citep{laya}, and the judging literature documents the failure modes here
precisely. MT-Bench and Chatbot Arena established the paradigm and their bias
taxonomy \citep{zheng-llmjudge}; relevance judges agree imperfectly with human
relevance \citep{sun-urara}; listwise re-rankers trade quality for instability
at scale \citep{rankgpt}; evaluators favour generations resembling their own
distribution \citep{panickssery-selfpref}; and verifiers can be gamed by
optimization pressure \citep{lightman-verify}. The contribution is not the
discovery of judge miscalibration, but the operational quantification of
same-domain inflation inside a deployed pipeline and the decomposition of which
judge \emph{role} is safe under which pool conditions
(Sections~\ref{sec:inflation} and \ref{sec:gate}).

These approaches verify \emph{during generation} with the generator itself.
This platform externalizes verification into a dedicated model that sees only
(query, document text) pairs and is consumed as a tool that any lane or agent
can call and whose verdicts fail closed, which makes verification an
inspectable, reusable service rather than a property of one generator's
prompting.

\emph{Abstention.} Beyond selective prediction under domain shift
\citep{kamath}, work probes whether LLMs know what they do not know
\citep{yin-knowwhat} and orchestrates models to abstain on knowledge gaps
\citep{feng-abstain}. Section~\ref{sec:abstention} locates the abstention of
this platform not in calibrated uncertainty but in the reading-and-answer
structure, a complementary architectural answer to the same problem.

\subsection{Agentic retrieval and tool protocols}
ReAct established interleaved reasoning over tools \citep{react-yao} and
tool-learning extended this to self-supervised tool use \citep{toolformer}; the
Model Context
Protocol standardizes tool servers as a machine interface \citep{mcp};
agentic-RAG surveys systematize the design space and list efficiency among open
challenges \citep{singh-agentic-rag}; Modular RAG frames retrieval as
reconfigurable modules \citep{modular-rag}. DeepRead shows that structure-aware
navigation over long documents improves agentic search \citep{deepread}, which
the librarian lane mirrors. The contribution is a policy layer: per-lane tool
whitelists, hard call budgets, and single-call server-side orchestration for the
hybrid lane (Section~\ref{sec:composition}).

\subsection{Positioning}
Table~\ref{tab:positioning} compares representative systems along the
dimensions a knowledge-base management system has to satisfy.

\begin{table}[!htbp]
\centering
\small
\setlength{\tabcolsep}{3pt}
\caption{Positioning against representative systems. \emph{Parse}: a quality
gate that rejects unusable documents before they enter the base. \emph{Org.\ \&
lifecycle}: organization and governance of the base beyond a flat index
(catalog, metadata and tag control, duplicate control, move/merge/delete).
\emph{Experience}: operational lessons kept as first-class retrievable
knowledge. \emph{Verify}: every cited item is judged against the query before it
may support an answer. \emph{Abstain}: explicit not-found is terminal, not a
fallback to parametric or web knowledge. \emph{Lanes}: several retrieval
strategies behind one interface under enforced budgets. \emph{Protocol}: a
standardized agent tool protocol as the machine interface. \yes{} present;
\no{} not addressed by the cited work; \textit{n/r} not reported in the
accessible source; italic entries name the mechanism where a bare mark would
mislead. Marks denote capability, not measured quality.}
\label{tab:positioning}
\begin{tabular}{@{}>{\raggedright\arraybackslash}p{4.0cm}ccccccc@{}}
\toprule
 & \rotatebox{90}{Parse} & \rotatebox{90}{Org.\ \& lifecycle} &
\rotatebox{90}{Experience} & \rotatebox{90}{Verify} &
\rotatebox{90}{Abstain} & \rotatebox{90}{Lanes} &
\rotatebox{90}{Protocol} \\
\midrule
Vanilla RAG \citep{rag-lewis} & \no & \no & \no & \no & \no & \no & \no \\
GraphRAG \citep{graphrag} & \no & \emph{communities} & \no & \no & \no & \no & \no \\
RAPTOR \citep{raptor} & \no & \emph{tree} & \no & \no & \no & \no & \no \\
LLM judge re-ranker \citep{rankgpt} & \no & \no & \no & \emph{ordering} & \no & \no & \no \\
Self-RAG \citep{selfrag} & \no & \no & \no & \emph{generator} & \emph{partial} & \no & \no \\
CRAG \citep{crag} & \no & \no & \no & \emph{corrective} & \no & \no & \no \\
DeepRead \citep{deepread} & \no & \emph{structure} & \no & \no & \no & \emph{agentic} & \textit{n/r} \\
ECM platforms \citep{paivarinta-ecm-2005,tyrvaainen-ecm-2006} & \textit{n/r} & \emph{lifecycle} & \textit{n/r} & \no & \no & \no & \textit{n/r} \\
Public KBs, DBpedia/YAGO \citep{auer-dbpedia-kb-2007,hoffart-yago2-kb-2013} & \textit{n/r} & \emph{ontology} & \no & \no & \no & \no & \no \\
\midrule
QDCVR (this work) & \yes & \yes & \yes & \yes & \yes & \yes & \yes \\
\bottomrule
\end{tabular}
\end{table}

Each tabulated system owns one or two dimensions: GraphRAG organizes a corpus as
a graph, RAPTOR as a tree, Self-RAG and CRAG verify with the generator, LLM
judges order candidates, ECM frameworks govern content across its lifecycle. No
tabulated system combines a quality-gated acquisition pipeline, a governed
lifecycle layer, first-class operational experience, per-citation verification,
enforced abstention, budgeted multi-lane retrieval, and a standard agent tool
protocol.

This paper contributes that combination as one design: contracts that hold
across acquisition, organization, and verification
(Sections~\ref{sec:acquisition}--\ref{sec:abstention}), not capabilities
implemented per component. It does not propose a better ranker. On the
gold-annotated corpus BM25 reaches Hit@5 of 1.00 while the judged hybrid reaches
0.44; the value of the gate measured here is auditable, fail-closed,
per-citation provenance, not ranking quality (Sections~\ref{sec:e4} and
\ref{sec:inflation}).

\section{System overview and design principles}
\label{sec:system}

QDCVR (v2.3) is a self-hosted knowledge-base management system that turns a
folder of PDFs, Office documents, and scans into knowledge bases that agents
can be trusted to answer from. It is deployed as four cooperating services
behind one configuration file: a Python service that schedules parsing and
owns the OCR engine, a TypeScript web application that serves the interface
and acts as the server-side proxy, an MCP server that exposes the platform's
capabilities as tools \cite{mcp}, and storage engines holding documents,
vectors, graph edges and task state.

\subsection{Requirements}
\label{sec:req}

The requirements below were derived from the deployment rather than
postulated, and each names the failure it prevents.

\begin{description}
\item[R1 --- Parse heterogeneous enterprise documents.] Process regulations,
failure reports, scanned inspection records and journal articles arrive as
PDFs, Office documents and images. The system must convert them to text it
can govern, including scanned pages, tables and formulas, and it must make a
parse failure visible rather than storing an empty or garbled document. The
100-paper acquisition left 21 papers that never became retrievable, and that
gap is reported as a first-class result with its cause left unattributed
(Section~\ref{sec:e1}).
\item[R2 --- Organize by content, not by folder.] A document's home must be
decided by what it says: library membership, tags, descriptions and graph
edges are derived from the content.
\item[R3 --- Manage the collection over its lifetime.] Libraries and
documents are created, moved, renamed, merged, split, deleted and re-indexed.
Every such operation must preserve agreement between the document on disk,
the tree registry and the library metadata, verifiably after the fact.
\item[R4 --- Retrieve with several independent channels.] Lexical, dense,
catalog and graph channels fail differently, so the system must generate
candidates with each, under explicit budgets, and expose which channel
proposed what.
\item[R5 --- Do not promote a similarity score into a citation.] A candidate
becomes an answer only after its fetched content has been judged to contain
direct evidence for the question. When no candidate survives, the system
returns an explicit not-found report instead of the closest text.
\item[R6 --- Be operable and engine-agnostic.] The platform runs on a single
local workstation (RTX 5060, 8\,GB, Windows), degrades rather than crashes
when a model or a subprocess fails, exposes long-running work as pollable
tasks, and drives whichever agent harness the operator has configured.
\end{description}

\subsection{Design principles}
\label{sec:principles}

Six principles shape the architecture, each enforced in the implementation
rather than left to prompt discipline.

\begin{description}
\item[P1 --- Content-routed organization.] Documents are organized by what
they say, not by where they sat: placement, tagging, description and graph
edges are derived from the parsed content, and retrieval starts from a
library catalog rather than a directory path (\S\ref{sec:organization}).
\item[P2 --- Content outranks similarity.] No similarity score promotes a
document into an answer. A judging model must read the candidate's fetched
content and vote, and the platform fetches that content itself so an agent
cannot argue for a document it has not read (\S\ref{sec:gate}).
\item[P3 --- Every module terminates at a contract.] Ingestion ends in a
checked post-condition, retrieval in a budgeted return value, verification in
a rubric score, and management in an audit.
\item[P4 --- Fail closed.] Missing text, an engine error, a timeout or an
unparseable verdict rejects the candidate rather than passing it, and a
document that cannot be retrieved after indexing is a failed ingest rather
than a warning.
\item[P5 --- Leave an audit trail.] Every candidate, verdict, cut decision,
tool call and lifecycle operation is recorded, so a claim about the system
can be re-examined from evidence instead of re-argued from memory.
\item[P6 --- Engine-agnostic agency.] Answer agents are pluggable harnesses
(14 adapters, including \texttt{omp}, Claude and ACP-compatible CLIs); the
platform constrains them through tool whitelists, budgets and prompts, never
by coupling to a particular vendor.
\end{description}

\subsection{Architecture and subsystems}
\label{sec:architecture}

Fig.~\ref{fig:arch} shows the four layers. The \emph{client and agent layer}
contains the web interface and the chat endpoint that adapts external agent
harnesses; requests are annotated with a lane policy before they leave this
layer. The \emph{tool layer} is an MCP server that exposes the platform's
capabilities as 95 tools in twelve groups (knowledge-base catalog, document
read/write, file system, parsing and task status, tags, search, the
verification gate, vector index and repair, knowledge graph, experience,
persona, and platform lifecycle and health);
the lane policy is a server-side whitelist, so a lane physically cannot call
a tool outside its set. The \emph{domain service layer} holds the tree and
catalog services, the parsing scheduler that manages the OCR engine as a
supervised subprocess, the vector and two-stage candidate services, the
experience service, the graph service and the resident judging worker. The
\emph{storage layer} holds the document tree with its registries, the vector
collections, the graph database and the task and session databases.

Three properties of this arrangement matter later. First,
the retrieval lanes of Section~\ref{sec:lanes} are not separate programs:
they are policies over one tool layer, which is why a new capability is
available to every lane at once and why lane behaviour is comparable across
harnesses. Second, long-running work (parsing, indexing, graph construction,
distillation) is exposed as a pollable task rather than a blocking call, so a
slow operation cannot stall an agent's turn. Third, the judging service is a
resident worker for engineering reasons, not architectural elegance: loading
the verifier per request cost 10--12\,s against 1--2\,s for a warm verdict,
which is why the worker, its file lock and its verdict cache exist
(\S\ref{sec:impl}).

\begin{figure}[t]
\centering
\includegraphics[width=\textwidth,height=0.86\textheight,keepaspectratio]{fig1_architecture}
\caption{QDCVR architecture. Documents enter through the acquisition
pipeline (left) and are parsed, quality-gated, split, described and indexed
into four stores. At run time (right), clients and agent harnesses reach the
platform through the chat endpoint, which attaches a lane policy enforced
against the MCP tool layer (95 tools); all retrieval lanes terminate at the
content-verification gate. Organization and life-cycle management (bottom)
keep the tree registry, the document files and the library metadata in
agreement.}
\label{fig:arch}
\end{figure}

Table~\ref{tab:modules} maps each function of the system to the components
that implement it, the design decision that shapes it, and the experiment
that provides evidence for it. The table is the paper's reading guide:
Sections~\ref{sec:acquisition}--\ref{sec:implementation} describe these rows
in order, and Section~\ref{sec:eval} evaluates them.

\begin{table}[!htbp]
\centering
\footnotesize
\setlength{\tabcolsep}{3.5pt}
\caption{System functions, their implementing components, the decision that
shapes each one, and the evidence that covers it. Evidence identifiers
E1--E8 refer to the evaluation subsections of Section~\ref{sec:eval}.}
\label{tab:modules}
\setlength{\tabcolsep}{5pt}
\resizebox{\textwidth}{!}{%
\begin{tabular}{p{2.5cm}p{3.9cm}p{4.5cm}p{1.8cm}}
\toprule
Function & Components and representative tools & Key design decision & Evidence \\
\midrule
Acquisition and parsing & parsing scheduler, OCR subprocess, duplicate
detector, split gate & quality-gated: a document that cannot be parsed is
rejected visibly, and a document that cannot be retrieved afterwards is a
failed ingest & E1 \\
Knowledge organization & content router, catalog, tag registry, description
generator & content decides placement; the catalog is the first retrieval
instrument & E1, E2 \\
Life cycle and integrity & tree registry, library metadata, re-index tasks,
three-source audit & re-indexing is explicit; every operation is audited
afterwards & E2 \\
Document graph & graph projection of documents, libraries and tags;
related-document queries & the graph is a navigation index, never an evidence
source & E2 \\
Catalog-first navigation & library and document listings, progressive
disclosure & metadata before bodies, so a candidate set is chosen without
loading content & E4, E5 \\
Content-verified retrieval & vector, lexical, catalog and graph channels;
lane policies; retrieval gate & verification is a promotion condition, not a
ranking hint; verdicts fail closed & E4, E6 \\
Operational knowledge & experience store, credibility tiers, distillation
runs & a lesson is a claim about practice, so it is drafted, gated and
capped until it is vetted & E3 \\
Agent interface and operations & MCP tool layer, chat endpoint, lane
whitelists, task registry & every capability is a tool with a schema; long
work is pollable; engines are interchangeable & E5, E7 \\
Persona and platform lifecycle & persona distillation and persona Q\&A tools;
version and update tools & shipped and reachable through the same tool
surface, but outside this paper's evaluation, so no claim is made here & -- \\
\bottomrule
\end{tabular}
}
\end{table}

\subsection{Data model and storage layout}
\label{sec:datamodel}

The platform keeps its data model readable: documents are
Markdown files, and the two registries are JSON and YAML files that a human
can open. Because they are files rather than a database, a document's address is a path and its identity a path plus a content hash, which is what makes
the three-source audit of \S\ref{sec:org-registries} possible. The vector
store is a separate engine that can disagree with the registries; the system
treats that disagreement as a detectable, repairable condition rather than as
an impossibility. Table~\ref{tab:datamodel} lists the stores and fields the
paper refers to.



\subsection{Object life cycle}
\label{sec:lifecycle}

A document moves through five states. It is \emph{raw} when it arrives, with
a content hash recorded. It is \emph{parsed} when text has been extracted and
has passed the parse-quality checks. It is \emph{stored} when the body, its
description and its tags have been written to disk and registered in both
registries. It is \emph{indexed} when the vector store and the retrieval
indexes contain it, which requires an explicit indexing call. It is
\emph{verified} when a retrieval probe returns the document for a query
derived from its own content. Only the last state counts as a successful
ingest; a document that stalls at \emph{stored} is an incomplete operation,
not a warning.

Libraries have a simpler life cycle: a library is \emph{active} while it is
addressable and searchable, and \emph{merged} or \emph{split} when an
operator consolidates or divides it. Merge and split are orchestrated
operator procedures over the document, catalog and index tools rather than
single tool calls, and neither was exercised under controlled conditions
(Section~\ref{sec:e2}); we describe the transition without claiming it.
Library metadata travels with the
library when it moves, including its distillation configuration, which is why
the configuration is stored per library rather than in the deployment file.
Platform-level versioning is explicit: the platform reports its version
(2.3.0 in this deployment, the value carried by the repository's
\protect\path{VERSION} file and returned by
\protect\path{ragctl version --json}) and can update itself from its
repository through
the same command-line interface that performs the other operations. The
update path is the one platform-lifecycle probe that failed in E2: its check
timed out after 30\,s, so version reporting is evidenced and self-update is
not.

Re-indexing is where the object life cycle meets retrieval quality. It is
explicit, asynchronous and observable: it returns a task identifier, it can
be forced, and its post-condition is a segment count that matches the
document content. Nothing in the write path re-indexes as a side effect.
This trades convenience for the property that no document can silently be
missing from the vector store: the cost is that a writer who forgets the
call leaves a stale index, and the mitigation is the audit of
\S\ref{sec:org-integrity}, which detects exactly that condition.

\section{Knowledge acquisition: parsing, ingestion, and index verification}
\label{sec:acquisition}

Acquisition is the stage this design treats most defensively; this
section reports the audit numbers we have and the measurement gaps we do
not.

\subsection{Multi-format intake and parsing}

Upload starts in the browser at \path{POST /api/parse/file-vt};
the Nuxt server forwards it to the asynchronous backend endpoint
\path{POST /api/v1/parse/file/vt}. The backend runs MinerU~\citep{mineru}
as a hidden subprocess on an ephemeral port, converting the source into
structured Markdown: OCR for scanned pages, layout-aware table recovery,
and formula extraction. The subprocess is reachable only through the
manager object, which resolves the assigned port at runtime; the service
never hardcodes it.
Its combined standard output and standard error go to
\path{backend/logs/mineru-api.log} and never to a pipe---an earlier
build failed on pipe closure with a bare I/O error, which appending to a
file removed. Parse jobs are polled with a 600\,s timeout, so a stuck
conversion ends in a recorded failure instead of blocking the caller.

The backend returns a \path{markdown_path}, a filesystem location
rather than document text; the proxy reads it back and writes the
document into the knowledge base through the tree filesystem service,
updating the on-disk Markdown, the tree index, and the per-library
metadata file in one operation. A path is transferred because parse
output can exceed a single JSON response and the proxy needs filesystem
access for extracted images.

\subsection{Quality gates before storage}

Every gate below runs before storage (Table~\ref{tab:acq-gates}): a
document that fails one is
refused, not stored and indexed anyway. Exact-duplicate detection hashes
the source with SHA-256 and compares it to the stored records before any
parse is scheduled; near-duplicates are detected in vector space with
\path{kb_find_duplicates(threshold=0.90)}; when the tool is
unavailable, a manual fallback channel matches filename and file size
and probes the vector store at 0.85. On the 322-part corpus the audit
recorded 0 true duplicates and 35 near-duplicate pairs, all adjacent
parts of one paper---the expected shape for a split document compared
against itself.

The parse-quality gate inspects the first 1{,}500 characters of the
markdown and refuses the parse on any of five signatures: OCR garbage
(more than 30\,\% garbled), binary residue such as null bytes or base64
fragments, headings without body text, whitespace collapse, or a
language mismatch between the source and the extracted text.

Storage has one rule: parsed documents are written with
\path{kb_doc_save_parsed}, never \path{kb_doc_create}. The
parsed path transfers the complete Markdown and copies all extracted
images; the generic create path truncates long content and drops images,
so using it on a parsed document silently loses both. Split parts go
through the create path individually because each is already a complete,
size-bounded unit.

\begin{table}[t]
\centering
\caption{Acquisition gates, their trigger, and what refusal means. Gate
names follow the packaged ingestion contract.}
\label{tab:acq-gates}
\small
\setlength{\tabcolsep}{3pt}
\begin{tabular}{@{}>{\raggedright\arraybackslash}p{2.1cm}>{\raggedright\arraybackslash}p{4.7cm}>{\raggedright\arraybackslash}p{5.4cm}@{}}
\toprule
\textbf{Gate} & \textbf{Trigger} & \textbf{On refusal} \\
\midrule
Exact duplicate & SHA-256 hit & Skip the source; report the existing path \\
Near duplicate & Vector similarity $\geq 0.90$; manual channel $\geq 0.85$ & Skip; supplement tags or description only \\
Parse quality & Garbled, empty, binary, or language-mismatched body & Reject the parse; re-parse or report the source \\
Split ceiling & Body beyond 30{,}000 characters & Split into parts; never store the oversized original \\
Atomic unit & Indivisible unit larger than the target & Keep whole with \path{oversized_atomic_unit}; hard cuts disabled \\
Tag anchoring & Tag absent from the read body sample & Return to extraction; do not release the document \\
Description readback & Claim not supported by the body & Rewrite the description; never the body \\
\midrule
Index verification & Collection wrong or \texttt{total\_chunks} $= 0$ & Ingest incomplete; do not report success \\
\bottomrule
\end{tabular}
\end{table}

\subsection{Structure-aware splitting and addressability}

Documents longer than a 30{,}000-character ceiling---a budget, not a
format limit, that keeps a stored unit inside the context window of the
agents that read it and bounds the per-query read cost of a single
hit---are not stored as single files. A splitting pass proposes
boundaries from the document's own structure, and a validator enforces
a target utilisation of 0.85, refusing oversized atomic units: a block
whose division would cut an equation or table stays whole. Each part is
written as its own file labelled \texttt{part $k$ of $N$} and read
through a line-offset address, so retrieval opens a section without
paying for the whole document.

The 100-paper arXiv corpus produced 322 stored parts, the corpus that
later audits and probes operate on. A separate 50-paper campaign under
the same 30{,}000-character setting produced 165 documents (nine whole,
156 parts) and 3{,}798{,}446 characters in 2{,}672\,s of parsing; the
same campaign at 2{,}000 characters produced 2{,}713 fragments, which is
what the ceiling prevents. The vector store chunks stored text at 500
characters with 50 of overlap~\citep{chroma,hnsw}; a comparison replica
index uses 800-character chunks and produced 4{,}528 chunks for the same
corpus.

\subsection{Content-derived descriptions and tags}

Descriptions and tags come from body text the ingestion agent has
actually read, not from filenames or directory positions. A description
must carry four elements---subject, method or technology, scenario or
problem, and key data or conclusion---with at least two concrete nouns,
and every claim is read back against the body before storage; a mismatch
rewrites the description, never the body. Tags pass a blacklist filter,
are normalised for case and language, and are re-checked against a
body sample: a tag absent from the text is dropped, not kept for
vocabulary reuse.

An earlier audit is why this gate exists: of 26 descriptions checked
against their own bodies, 24 did not match---the boilerplate
``Gutenberg front/back matter'' reached real chapters across a batch,
and the librarian that reads part labels selected two parts instead of
seven. The wider catalog carries the same defect, with 156 of 322
documents flagged for a broken part prefix. Description quality is
therefore an empirical result, not an assumption: descriptions are
claims about content, not facts about it, and a claim that can be wrong
at ingest scale must be cheap to verify and cheap to repair---verification
later became a platform primitive (Section~\ref{sec:lanes}), and the
acquisition audit ends with descriptions upgraded 165/165 rather than a
generation count.

\subsection{One ingest, end to end}
\label{sec:acq-trace}

The stages above compose into one path. An MCP-layer caller passes the
source path to \path{parse_doc}, which returns a task identifier at once;
\path{parse_task_status} is polled until the task is done. A task that
finished without a successful parse result stops the ingest there rather
than being saved as an empty document. The parse yields a
\path{markdown_path}; the proxy reads that file and writes it into the
tree registry, updating on-disk Markdown, tree index and library metadata
together. The parse-quality and duplicate gates run before anything is
stored, so a refusal ends the ingest. A part beyond the
30{,}000-character ceiling is split, and each validated part is then written as
a document of its own through \path{kb_doc_create}, addressed
\texttt{part $k$ of $N$} and carrying its own description, while the temporary
unsplit Markdown is deleted. Description and tag generation read the body, and
each description claim and each tag is checked back against a body sample
before it is stored. An unsplit document goes in through
\path{kb_doc_save_parsed}, which writes body and description together and
copies the extracted images; tags are attached afterwards by their own call.
Indexing is a separate explicit call---storing a document builds neither the
vector collection nor the graph---and the ingest ends with a retrieval probe
whose query is derived from the document's own content, the step at which the
document becomes \emph{verified} (Section~\ref{sec:lifecycle}). The per-paper
ingest checks and the ingest state file record what the gates observed for the
papers that reached them.

\subsection{Index construction and retrievability verification}

Indexing writes to three channels: a per-library vector collection, a
lexical BM25 index~\citep{bm25} over chunked documents, and a document
graph in Neo4j~\citep{neo4j}. Ingestion success is judged by observable
effects, not the tool's success envelope: a path that reported success
returned an empty index payload after a GPU out-of-memory failure,
so a document counts as indexed only with a non-empty collection name,
the expected collection identifier, and \texttt{total\_chunks} $\geq 1$.
Graph construction is asynchronous, and task status is not completion
evidence---an earlier build left a task row
stuck in \texttt{running} while its artifacts were queryable. The step
closes with an eight-point final check: the stored copy read back at 500
characters, the chunk and graph node counts, and three-layer metadata
consistency across the tree index, the library metadata file, and the
files on disk.

\subsection{Asynchronous execution, failure recovery, and observability}

Parsing and graph construction are non-blocking: each returns a task
identifier that a registry resolves across tool processes, so callers
poll status instead of holding a connection open. Batch ingestion is
resumable and was resumed in practice: restarting the client after an
interruption in a 100-paper run skipped documents already present and
continued from the recorded position. Per-paper parse time in the
50-paper campaign ranged from 16.3\,s to 40.4\,s over 26 subject
domains, 1{,}115.5\,s in total, an average of 22.3\,s per paper.

That bookkeeping has a limit. Of 100 submitted papers, the state file
records 91 names, 79 were retrievable in the knowledge base at
evaluation time, and 21 never became retrievable. The reported failure
mode is formula-dense scanned documents, and a mid-run client crash was
recovered by restart, but no artifact counts parse attempts, parse
failures, or the error text of a failed document: the 91-versus-79
difference is an observability gap we can report from the artifacts but
not attribute from them.

Cross-lingual ingestion shows the same gates holding without
language-specific tuning. The corpus is 400 Wikipedia pages per language
in English, Chinese, and Japanese, ingested in two rounds; the second
round was bit-identical to the first (Table~\ref{tab:acq-crosslang})
and ingestion success was 1.0 for all three languages in both rounds.
Membership---where each page landed---is 0.875 for English, 1.0 for
Chinese, and 1.0 for Japanese; storage completeness---split parts summed
back to the source length---is 0.975, 0.9559, and 0.9402; and
self-retrieval Hit@1, the score for the page's own indexed text, is
0.975, 1.0, and 1.0.



The 100-paper pipeline audit records what each gate observed---every paper
parsed and routed by content, nine per-paper ingest checks 100/100, content and
storage audits 322/322 and 33/33---but observability and retrievability are
different ledgers: 91 papers reached the state file and 79 were retrievable at
evaluation time (Table~\ref{tab:acq-pipeline}). What the gates do
not buy is accounting: with no counter of attempted or failed
acquisitions, the difference between 91 recorded and 79 retrievable
papers cannot be split into parse losses and lost records. We treat a
submission as knowledge only once verification has found it.

\section{Knowledge organization and lifecycle management}
\label{sec:organization}

Acquisition decides what enters the system; organization decides what the
system is. Here organization is a governed process rather than a folder
convention or an emergent property of embedding similarity, with three
invariants: content decides where a document belongs, the disk, the tree
registry and the library metadata must agree about every document, and every
operation that removes or relocates knowledge is preceded by a survey and
followed by a verification. This chapter describes their implementation and
where the three stores can still diverge.

\subsection{Content-routed placement}
\label{sec:org-placement}

Placement reads the parsed document, not its filename or its folder of
origin. The pipeline selects the library whose stated scope covers the
content, or proposes a new library when no existing scope does, writes the
document, derives a description and a tag set, and runs the same final check
used by acquisition before the document counts as placed. Two rules keep the
catalog usable as a retrieval instrument: an existing library is preferred
over a new one; and a document is placed once, because re-routing is an
explicit lifecycle operation rather than a side effect of a later ingest.

The decision itself is made over library scope rather than folder names. The
ingest agent is given the candidate libraries with their scope descriptions and
the parsed body of the document, and selects the library whose scope covers the
content. That assignment is recorded in the catalog and the document is written
to that library. The decision is per document, not per part: a split document's
parts inherit the assignment made for their source. When no scope fits, the
agent defines a new library and states its scope before writing into it. The one
ingest run we instrumented carried its assignments in a classification file
keyed by paper identifier and built from their content together with the library
descriptions; an unassigned paper raised an error rather than a guess. The
prototype fixes neither a model nor a similarity threshold for the choice, so it
remains operator-visible at review time.

The library tree supports one level of content grouping in which child
libraries are addressable through a parent. The implementation carries a
caveat we report rather than hide: the vector segments of child documents are
stored in the parent library's collection, so a search that targets a child
library identifier returns no candidates. The shipped operational guidance is
to search the parent collection and to use the child libraries for navigation
and scoping rather than as search endpoints. This is a limitation of the
storage mapping, not of the routing rule.

\subsection{Catalog and metadata registries}
\label{sec:org-registries}

Table~\ref{tab:datamodel} lists the fields of the two registries that
describe the same collection: the global tree (\texttt{.tree-fs.json}) is the
authoritative structure, and each library carries a
\texttt{.knowledge-base.yml} indexing its Markdown documents and their image
directory.

Agents read the catalog through a progressive-disclosure interface: the
listing returns names, descriptions and counts, document listings return
metadata without bodies, and only an explicit read returns content,
attributed to a specific document. A caller can therefore narrow a search to
a handful of candidates without loading a body, which makes the catalog lane
viable at 5--6 turns and 54--91~s (Table~\ref{tab:lanes}).

Correctness of the catalog is three-source consistency: for every document,
the file on disk, its entry in the tree registry, and its record in the
library metadata must agree on existence, path, size and index state. The
platform verifies this agreement rather than assuming it: in the acceptance
campaign a renamed document kept its old vector segments, so the recorded
chunk count read 18 where 9 chunks existed, and the document matched its own
former path at 0.9738 similarity, high enough to be reported as a
near-duplicate pair. Only an explicit force re-index repaired it
(18~$\rightarrow$~9). Metadata and index consistency are therefore two
different obligations: the first is satisfied by a single write, the second
is not, because the vector store is updated by a separate call.

\subsection{Tag taxonomy and duplicate detection}
\label{sec:org-tags}
\label{sec:org-dedup}

Tags are the coarse, human-legible axis of the collection, and retrieval by
tag is a first-class operation. A tag registry holds the vocabulary; each
document's tags are normalized, filtered against a blacklist of structural
strings (section headings, generic words such as ``introduction''), and
re-checked against the document body so that a tag must occur in the content
it labels. Governance is a maintenance loop rather than a schema: no taxonomy
is imposed on authors and the vocabulary grows with the corpus, so the audit
reports tags with zero referencing documents, accumulated from test documents
and chapter headings, and removes them only when instructed. In the audited
campaign 165 documents passed the tag gate with no failures, and the live
registry held 316 tags. We do not claim a controlled vocabulary; we claim
that the vocabulary is inspectable and that orphaned entries are detectable.

Deduplication is a collection-level question as well: the same content can
arrive through two ingest paths or into two libraries. The procedure is
survey-then-act --- enumerate candidate duplicate groups at a similarity
threshold (0.90, with a manual fallback channel at 0.85), inspect each group,
delete nothing automatically --- and the audited run illustrates why the
survey step matters more than the threshold. On the 322-document corpus the
detector reported 35 candidate groups with similarities between 0.90 and
0.975, every one of them a pair of adjacent parts of the same long paper, and
no deletion was performed; a policy that deleted on detection would have
destroyed 35 legitimate document parts. The same run confirmed zero true
duplicates. Cross-library duplication is detected but not merged
automatically, because merging two libraries carries provenance consequences
and stays with the operator.

\subsection{Library lifecycle operations}
\label{sec:org-lifecycle}

Table~\ref{tab:org-lifecycle} lists the operations the platform exposes, the
interface that performs them, the consistency obligation each one carries,
and the state observed for each.

\begin{table}[t]
\centering
\small
\caption{Library and document lifecycle operations, their consistency
obligation, and the state observed in the acceptance campaign. ``Audited''
means the three stores were compared after the operation.}
\label{tab:org-lifecycle}
\begin{tabular}{@{}p{0.22\textwidth}p{0.24\textwidth}p{0.20\textwidth}p{0.24\textwidth}@{}}
\toprule
Operation & Interface & Obligation & Observed \\
\midrule
Create library & library create & tree + metadata + directory & audited \\
Rename or move document & document move & path updated in all three stores & audited; stale segments until force re-index \\
Update description or tags & document update & metadata rewritten, body re-checked & audited (165/165) \\
Update content & document update & new hash, re-index required & re-index not implicit \\
Delete document & document delete & removal from all three stores & audited (1~$\rightarrow$~0) \\
Batch delete & batch delete & per-document verification & 1/1 in campaign \\
Move library & library move & subtree and metadata move together & audited (0~$\rightarrow$~1) \\
Merge libraries & library merge & provenance preserved & demonstrated, not controlled \\
Split library & library split & no part orphaned & analysed, deliberately not executed \\
Rebuild index & re-index (task id) & segment count matches content & 18~$\rightarrow$~9 after force \\
\bottomrule
\end{tabular}
\end{table}

The rule the table carries is that re-indexing is not implicit: creating,
moving or rewriting a document does not update the vector store, only an
explicit index operation does, and the audit treats a missing or stale index
as an incomplete operation rather than as a warning. The contribution here is
a contract under which disagreement between an approximate index and a
mutating document store becomes detectable, not a new maintenance algorithm
\citep{xu-spfresh-2023}. The evidence boundary is narrower: move, rename,
update, delete and re-index were exercised with real tool calls and audited
afterwards; merge and split are implemented and were demonstrated on live
libraries, but they have no controlled evaluation, and the sub-library
hierarchy was analysed and deliberately not executed pending an operator
decision. No concurrent organizational operation was tested.

\subsection{The document graph}
\label{sec:org-graph}

The platform maintains a property graph with document, library, tag, folder
and entity nodes connected by membership, tagging, relatedness, identity and
containment edges. Those edges derive from tag co-membership, vector
similarity above a threshold, and the verdicts produced by content
verification during retrieval. Adding one document to a library added four
edges in the audited run (4{,}675~$\rightarrow$~4{,}679), one per assigned
tag; a full build over the 100-paper corpus produced 658 nodes covering 333
documents and 316 tags, and the 165-document intermediate state held 335
nodes and 1{,}671 edges.

The graph keeps a deliberately restricted role. It is a navigation instrument
that expands a candidate set and surfaces documents sharing context across
libraries. It is never evidence: as a standalone retriever it reached Hit@5
of 0.5625 on the industrial corpus and 0.2 on the second corpus, well below
the lexical and dense channels (Table~\ref{tab:retrieval}). Graph-based
organization in this system is useful for discovery and weak as a ranker.

Two implementation gaps belong here. Graph construction is not yet guaranteed
by the ingestion path: a document ingested through one agent harness received
no graph node, so a library reported two of three documents. And batch graph
builds have been observed to leave a task record in the running state after
the work completed, so task status cannot be used as completion evidence; the
graph itself must be queried.

\subsection{Operational knowledge as a first-class collection}
\label{sec:org-experience}

Process documents describe what the organization intends; operational lessons
describe what actually happened. The platform stores the second kind as a
first-class collection with the same addressing, indexing and retrieval
interface as ordinary documents, but with two differences that reflect the
epistemic status of a lesson. First, an experience carries a credibility tier
that controls whether it may be shown (Table~\ref{tab:exp-tiers}); the
experience channel's similarity threshold is 0.55 against the document
channel's 0.35, and an unvetted or disputed entry is capped even when it
scores well. Second, entries are drafted before they are served: extraction
produces candidates from ingested content or from recurring question-answer
pairs, and automated distillation runs as a separate agent process on a
schedule, guarded by a circuit breaker (three consecutive failures suspend it
for 24 hours) and a heuristic fallback marked low-confidence. Nothing reaches
the served collection without an explicit approval step, the same fail-closed
posture the retrieval gate applies to citations.

\begin{table}[t]
\centering
\small
\caption{Experience credibility tiers and modifiers. Tiers combine retrieval
score, content-verification score, ratings and review count; modifiers cap
the tier rather than promote it.}
\label{tab:exp-tiers}
\begin{tabular}{@{}l>{\raggedright\arraybackslash}p{3.6cm}p{0.46\textwidth}@{}}
\toprule
Tier & Condition & Effect \\
\midrule
P0 strong & sim $\ge 0.65$ $\wedge$ content $\ge 6$ $\wedge$ rating $\ge 4$ $\wedge$ reviews $\ge 1$ & may be presented first \\
P1 reference & sim $\ge 0.45$ $\wedge$ content $\ge 4$ & presented with credibility noted \\
P2 weak & sim $\ge 0.35$ $\wedge$ content $\ge 3$ & suppressed unless expanded \\
discard & content verification fails or sim $< 0.35$ & never presented \\
\addlinespace
disputed & reviews $\ge 3$ $\wedge$ rating $< 2.0$ & capped at P2 \\
unvetted & 0 reviews $\wedge$ 0 applications & capped at P1 \\
auto-extracted & not human-reviewed & capped at P1 until applied twice and rated $\ge 4$ \\
\bottomrule
\end{tabular}
\end{table}

The evidence for this subsystem is deliberately case evidence. One extraction
campaign produced thirteen candidate drafts; one end-to-end scenario
sedimented a single experience, which a later global experience search
returned first with a local similarity of 0.711; and the approval path was
never exercised, because the draft pool was empty when it was tested. A
single scenario and thirteen candidates cannot support an effectiveness
claim, and we do not make one. What the campaign establishes is that
operational knowledge travels through the same pipeline as documents,
addressed, indexed, gated and auditable, rather than living in a wiki or a
prompt.

\subsection{Integrity verification as routine}
\label{sec:org-integrity}

Every property claimed in this chapter is checked by a verification pass
against the live collection, so that consistency, integrity and health are
measured rather than asserted. The pass is read-only by default; repairs
require an explicit instruction, and the report is produced first, because an
audit that repairs as it goes cannot report what it found.

The pass is not ceremonial. It is how the two index-integrity incidents of
this deployment were found, a silently lost collection and a stale index, and
how the reorganization campaign was certified after the fact: 322/322
metadata checks, 165/165 description checks, 33/33 sampled consistency
checks, and 15/15 retrieval regression probes with no degradation. In the
same campaign the audit surfaced the failures reported above: the rename-time
chunk duplication, the missing graph node, and a task-status record that
outlived its task. A collection whose health is a command is a collection
whose health can be argued about with evidence.

\subsection{What organization buys, and what it does not}
\label{sec:org-summary}

Organization buys candidate quality and auditability. Because the catalog is
curated and consistent, an agent can select a library and a document set
without loading bodies, which is the mechanism behind the catalog lane and
behind cross-library discovery in an environment with nineteen libraries.
Because every placement, tag, description and deletion left a record, the
collection can be re-verified after the fact; the numbers in this chapter are
audit outputs rather than survey estimates.

Three things organization does not do, and the paper does not claim. It does
not repair a bad parse: a document whose parse lost its tables cannot be
recovered by placing it well. It does not prevent divergence by construction:
the write path can leave the vector index behind, and only an audit and an
explicit re-index restore agreement. And it does not scale without limits:
the largest collection we report is 322 documents in five libraries, and the
merge and split operations are unevaluated. These are engineering boundaries.

\section{Content-verified retrieval}
\label{sec:retrieval}
\label{sec:lanes}

Every citation a lane proposes must pass a content-verification gate before it
may appear in an answer. The gate scores each retrieved segment on a 0--8 rubric
(topic 0--3, scenario 0--3, evidence 0--2) and keeps a document only if at least
one of its segments reaches the acceptance threshold on that rubric's $[0,1]$
projection: 0.5 in general, 0.3 for enumeration questions, corresponding to the
rubric's direct-answer band (Section~\ref{sec:gate}). It processes at most six
references per call, takes the
matched text from the server side rather than from the caller, and fails closed.
No surviving evidence means the caller must answer that the corpus does not
contain the answer.

\subsection{One query through the gate}
\label{sec:trace}

Q2 of Section~\ref{sec:eval-setup} is \emph{``EL inspection shows edge delamination:
root cause and corrective actions''}, asked through the public chat API with the
failure-analysis library pinned and the Hybrid lane selected.

\begin{enumerate}
\item \textbf{Policy.} The chat endpoint attaches the Hybrid lane's tool
whitelist server-side, so the agent cannot call a tool outside it.
\item \textbf{Candidate generation (one call).} The agent issues a single
retrieval call. In the shipped tool layer that call is a server-side
composition of \texttt{kb\_search\_vector}, the catalog tools \texttt{kb\_list}
and \texttt{kb\_get\_documents}, and the judging tool \texttt{kb\_laya\_judge};
\texttt{kb\_hybrid\_search} was an earlier-generation wrapper that is no longer
registered, so the configuration evaluated below was reproduced by composing
those tools. Two server-side channels then run: a dense vector channel
(the \texttt{kb\_search\_vector} call, top-30, cosine threshold 0.35) and a
catalog keyword channel over the pinned library's descriptions. Candidates
merge on (library, document), agreed documents rank first, and the pool is
capped at twelve.
\item \textbf{The gate.} The server fetches each candidate's text, segments it,
and scores \emph{does this text contain direct evidence that answers the query}
in batches of at most six. The agent supplies no text, so it cannot
argue for a document it has not read. On Q2 the gold report scores 0.846 while
five extruder-telemetry reports score 0.887--0.926.
\item \textbf{Cut and read budget.} A document-level relative cut keeps scores
$\geq \mathit{best}-0.15$ with a top-5 floor and a cross-domain anchor guard, so
an out-of-domain score cannot set the cut line; the top five survivors return in
full, at most 2\,000 characters each. That read window is the agent's only
full-content channel; the rest of the pool survives as scores in the audit
trail.
\item \textbf{The answer.} Citations name the document path, the lane, and the
engine score, graded P0 (direct evidence) through P2 (weak).
\item \textbf{The terminal paths.} If the cut leaves no survivor, Lane~A
escalates once to Lane~B; if Lane~B exhausts its budget, the run reports
not-found.
\end{enumerate}

Where the design fails, the same trace is explicit. Under the shipped re-ranking
role, Q2's gold (0.846) falls to rank eight behind the five higher-scoring
telemetry reports, outside the five-document read window, so the agent reads
confusers (Fig.~\ref{fig:anatomy}). Under the filter role the same pool, scores,
and verdicts leave the gold first: same gate, same evidence, only the role of the
induced ordering changes. This foreshadows the central result
(Section~\ref{sec:inflation}).

\subsection{The gate contract}
\label{sec:gate}

The gate is a locally hosted evidence-classification model, the Laya judging
service \cite{laya}, invoked through a \texttt{kb\_laya\_judge} tool. Given a
query and a batch of references, the server fetches and segments each document
and scores every segment under one of three criteria: \emph{evidence} (specific
substantiation; a topic match alone is insufficient), \emph{instance} (the text
contains a member of the requested class, routed automatically for enumeration
questions), or \emph{auto}. A document survives only if its best segment reaches
the threshold, so one at vector score 0.40 with a \emph{yes} verdict survives
while one at 0.95 that the judge rejects is dropped.

The score is documented as a 0--8 rubric (topic 0--3, scenario 0--3,
evidence 0--2; $\geq$6 answers directly, 5 escalates to the librarian lane,
$\leq$4 drops), with the $[0,1]$ score as its machine-readable projection. That
projection saturates (the logged maximum score is 1.0000), one reason thresholds
discriminate poorly between golds and same-domain confusers
(Section~\ref{sec:inflation}).

Every operating constant is a shipped production default: the thresholds, the
six-reference judge batch, the hybrid candidate cap of twelve, and the 0.15
relative-cut margin with its top-5 floor. They were fixed during earlier
deployments on a different corpus and were not tuned on the evaluation questions
of Section~\ref{sec:eval-setup}. The paper states their status rather than
selling the operating point: on the evaluated corpus none of the three
thresholds ever fired, because every candidate the retrievers returned already
scored far above them (Section~\ref{sec:inflation}). The gate's value there is
fail-closed provenance and the filter role, not discrimination; a deployment on
other content would have to re-tune these constants, and Section~\ref{sec:threats}
records that as a threat rather than a result. The classifier is a single, locally fine-tuned
evidence classifier whose checkpoint reference ships with the platform
artifacts; it is deterministic in our replays, and every verifier finding in this
paper inherits that choice, a scope limit noted in Section~\ref{sec:threats}.

Fail-closed semantics apply to every failure mode: fetch errors, engine errors,
and timeouts count as rejection. Survivors aggregate into an evidence pack with
per-passage provenance.

\subsection{Three lanes}

The three lanes are three points on a latency/coverage frontier
(Table~\ref{tab:lanes}; Fig.~\ref{fig:lanes}), not three implementations of one
idea.

\begin{table}[!htbp]
\centering
\scriptsize
\setlength{\tabcolsep}{3.5pt}
\caption{The three lanes as operating points. Measured figures are the
round-1 agentic runs of Section~\ref{sec:agentic}. All three share one
verification gate and one honesty contract; they differ only in retrieval
affordance and budget.}
\label{tab:lanes}
\begin{tabular}{p{1.85cm}>{\raggedright\arraybackslash}p{3.05cm}>{\raggedright\arraybackslash}p{2.15cm}>{\raggedright\arraybackslash}p{2.25cm}>{\raggedright\arraybackslash}p{3.15cm}}
\toprule
Lane & Optimises for & Tool budget & Measured (round 1) & Use when \\
\midrule
A. Search & latency on pinned libraries & 20-turn cap & 6--8 turns; 70--165\,s & interactive default; single-library questions \\
B. Librarian & exhaustive recall across all libraries & hard 12-call budget & 5--6 turns; 54--91\,s & must-not-miss; enumeration questions; no library pinned \\
C. Hybrid & single-call server-side orchestration & exactly 1 call & 3--5 turns; 36--70\,s & fastest grounded answer; latency-critical \\
\bottomrule
\end{tabular}
\end{table}

\begin{figure}[t]
\centering
\includegraphics[width=\textwidth]{fig2_lanes}
\caption{The three retrieval lanes. Every lane ends at the same
content-verification gate and the same honesty contract: zero verified
evidence escalates (Lane A$\to$B) or terminates in an explicit not-found
report.}
\label{fig:lanes}
\end{figure}

\textbf{Lane A (Search).} Phase A0 rewrites the query into (subject, attribute,
constraints) form, splits comparison questions into parallel sub-queries, and
routes incident-style queries to the experience store
(Section~\ref{sec:org-experience}) first. Phase A1 issues one wide dense query (top-30,
threshold 0.35, cross-library balancing), deduplicates by document, and keeps
candidate documents rather than truncating to top chunks. Phase A2 submits
citation-only references (library id plus document path) to the gate, and Phase
A3 returns the judge's evidence pack, this lane's only full-content channel, from
which the agent composes an answer whose citations carry confidence levels
(P0 direct evidence, P1 partial, P2 weak).

\textbf{Lane B (Librarian).} A fixed spine under a hard budget of twelve tool
calls. L0 reads the library catalog (ids, names, descriptions, counts) and prunes
shelves by the rewritten query; L1 lists each shelf's documents with
descriptions; L2 routes by query shape, sending literal and enumeration questions
(``which documents mention X'') to a deterministic scanner
(\texttt{scripts/enumerate\_scan.py}, which pulls the shelf's documents through
the tool layer and matches the requested terms literally and by regular
expression: $\approx$0.5\,s per document, zero LLM turns) and semantic questions
to judging; L3 judges in batches of at most six,
because larger batches historically triggered transport timeouts
and fail-closed batch loss; L4 reads only survivors, bounded to 2\,000 characters
and starting at the judge-reported evidence line; L5 emits a five-part report
(search paths, answer, sources, confidence, blind spots). The query-shape
choice is measured, not stylistic: on a probe with nine documented mentions, the
semantic judge recalled three of nine under the instance criterion at both
thresholds, while the scanner recalled nine of nine in 11.3\,s. Two rules are absolute:
catalog descriptions never prune documents by themselves, and an exhausted budget
terminates in blind-spot disclosure, not in a guessed answer.

\textbf{Lane C (Hybrid).} Agentic latency is dominated by LLM turns, not tools
(Section~\ref{sec:composition}), so Lane C spends exactly one call by moving the
pipeline into the server. A single call executes P1 (a
parallel block of a vector channel, top-30 at threshold 0.35, and a catalog
keyword channel); in the shipped tool layer that call is the composition of
\texttt{kb\_search\_vector}, \texttt{kb\_list}, \texttt{kb\_get\_documents} and
\texttt{kb\_laya\_judge} described above, not a single wrapper tool. P3 merges
candidates by (library, document), agreed-by-both
documents first, capped at twelve; P4 judges every candidate with the gate
(criterion \emph{evidence}, threshold 0.5); P4b applies the document-level
relative cut; and P5 reads the top five documents in full ($\leq$2\,000
characters each). The complete audit trail ships in the response, so a one-call
lane is still an inspectable lane.

\subsection{Selecting a lane is a budget decision}

The three lanes trade latency for recall under an enforced call and turn budget,
so the choice is how much budget to spend before the gate runs, not which
retriever is better.

What makes the lanes interchangeable is the gate. All three terminate in the same
contract: one gate, one fail-closed semantics, one evidence pack. Replacing a
lane changes the retrieval affordance and the budget, not the meaning of a
citation. Escalation follows a trigger, not a hunch: zero verified evidence
escalates Lane~A to Lane~B exactly once. Query-shape auto-selection is future
work (Section~\ref{sec:threats}).

One caveat is honest to record. On the evaluated corpus Lane~A occupies no
\emph{distinct} operating point, because Hybrid is faster and Librarian more
exhaustive. A two-lane deployment (Hybrid plus Librarian) is therefore defensible
today, with Search retained as the pinned-library interactive default; because
the lanes share one gate and one abstention semantics, that reduction costs
coverage, not trust.

\subsection{The two-stage entry point}
\label{sec:twostage}

One configuration scored in Section~\ref{sec:e4} is not one of the three lanes:
\path{kb_search_two_stage(query, kb_id, stage1_top_k=20, stage2_top_k=5, enable_graph_expansion=True, score_threshold=0.0, balance_kbs=False, slim=False)}.
It serves callers that want recall without paying for a verification lane, and
the platform's own retrieval skill names it the preferred entry point for agent
searches. Stage~1 builds a candidate document pool from the lexical index and,
when graph expansion is on, from the graph neighbours of the matched nodes
($\leq$20 documents); stage~2 runs the dense retriever inside that pool and
returns up to five chunks per document. A cross-library search that returns
documents from fewer than two libraries triggers an automatic supplementary
vector pass---the same lexical blind spot the librarian lane corrects by
escalating. It is fast because it never calls a judge: Hit@5 is 1.00 at 77\,ms
mean (Table~\ref{tab:retrieval}). That is also its limit, and the reason it is
not a lane: stage~2 returns chunks, so a caller that stops here holds candidate
text with no verdict attached, and any claim built on it is unverified by
construction. Only the gate turns a candidate into evidence.

\subsection{The gate as a platform primitive}
\label{sec:gate-primitive}

The gate is a server-side tool rather than a prompt:
\path{kb_laya_judge(query, documents, criterion=auto, threshold=0.5, max_evidence_chars=40000)}. The platform, not the agent, chooses the criterion
and fixes the threshold. Calls are serialized through a single judge worker with
a 900\,s call timeout, because an early concurrent workload crashed the worker
with a torch paging-file error; every verdict is logged with its query,
references, scores, and decision.

That placement makes verification a layer property rather than an agent
aspiration: any lane, harness, or external client inherits the gate, and none can
skip or re-threshold it per query. Enforcement has a second half in the lane
policy, a server-side per-lane tool whitelist: a librarian-lane agent is denied
the vector and two-stage tools and cannot reach around the catalog-first path it
was given. Denial is a property of that lane, not of the platform: the same
two-stage tool is scored directly as its own configuration in
Section~\ref{sec:e4}, and the vector tool backs Lane~A.

Externalized verification also distinguishes this design from verifiers inside
the generator. Self-RAG and CRAG verify during generation with the generator
itself \cite{selfrag,crag}; the gate sees only (query, document text) pairs and
is consumed as a tool. It is therefore reusable and inspectable, and equally
fallible: the LLM-judge critique applies here as much as anywhere
\cite{zheng-llmjudge}, which is why verdicts fail closed, are logged, and are
reported as evidence rather than as truth (Section~\ref{sec:inflation}). Deployment
cost is paid once: a twelve-candidate verification adds $\approx$0.5\,s to a call
(Table~\ref{tab:retrieval}), and a cold judge load costs 10--12\,s
(Section~\ref{sec:impl}).

\subsection{Honest abstention}

A threshold that never rejects a candidate cannot refuse a query. On four
out-of-corpus probes the judged hybrid still keeps 9--11 of its 12 candidates
(Section~\ref{sec:abstention}, Fig.~\ref{fig:ood}), and the graded probes behave
no differently: prompts with partial coverage and prompts with none both keep
10--12 of 12 candidates, with top scores up to 0.983
(Table~\ref{tab:robustness} and Section~\ref{sec:e7}). Because the $[0,1]$
projection saturates, no cut point separates a supported answer from an
unsupported one.

Abstention therefore lives in the reading-and-answering layer, enforced by the
lane contract at answer time rather than advised in a prompt. A caller with no
surviving evidence must report a not-found with its search scope: which libraries
and how many documents were searched, how the query was routed, and which topics
sit outside the corpus. The escalation rule makes that state terminal, so an
exhausted budget produces a report and not a guess.

Abstention is thus a terminal state evaluated like any other outcome: all six
out-of-corpus probe runs returned explicit not-found reports with zero fabricated
citations, even though the tool layer had handed them candidates
(Sections~\ref{sec:agentic} and \ref{sec:abstention}). The division of labour is
the point: the gate decides what may be cited, the reading layer decides whether
to answer at all, and the measurements put the reliability in the second and not
the first.

\section{Agent interface, implementation, and non-functional properties}
\label{sec:implementation}
\label{sec:impl}

\subsection{The tool surface}
\label{sec:impl-tools}

The platform is exposed to agents as a single Model Context Protocol
server~\cite{mcp} carrying 95 tools in twelve groups: knowledge-base catalog;
document read/write; file system; parsing and task status; tags; search; the
verification gate; vector index and repair; knowledge graph; experience;
persona; and platform lifecycle and health. The count is the number of registered tool decorators in
\texttt{kb-mcp/server.py}, checked one-to-one against the function
definitions they register; the project's own documentation had drifted to a
lower number, so we report the measured one. Of the twelve groups, the paper
evaluates seven (catalog, documents, parsing, search, gate, index and graph);
the persona subsystem and the platform-lifecycle tools are outside its scope,
and we say so rather than claiming coverage of the whole surface. The interface is not a reduced
agent-only view: it owns no private data path. Writes go through the same
HTTP services the user interface calls, reads read the same
\texttt{.tree-fs.json} and \texttt{.knowledge-base.yml} registries the web
server reads, and the MCP process itself contains no HTTP code, only a
client module.

Three governance rules make a surface this wide safe to hand to a model.
First, read-path tools and write-path tools are disjoint families, and
retrieval lanes are granted the read family only. Second, a stored parsed
document cannot enter a knowledge base through the generic document writer:
it must pass the parsed-document writer (\texttt{kb\_doc\_save\_parsed}),
the single entry point that carries parse provenance and forces the
ingestion quality gates to run. Third, long-running operations---a parse
batch, a re-index, a graph build, a meditation pass (the scheduled pass that
rewrites the experience store's rolling summary from accumulated questions,
retrievals and feedback)---return a task id
immediately instead of blocking the session, and are polled through the
task-status tools backed by an in-process task registry.

\subsection{Harness integration and lane policy}
\label{sec:impl-harness}

One chat endpoint adapts fourteen agent harnesses to a uniform message
stream, so a vendor CLI, an ACP-compatible agent, and the platform's own
engine all reach the same retrieval lanes. The lane policy of
Section~\ref{sec:lanes} is enforced here rather than by prompt discipline:
each lane is attached to a tool allow-list, of which 21 entries are
read-only MCP retrieval tools, and to a deny-list of more than 40
harness-internal tools (file editors, web search, sub-agent spawning). The
gate was observed working in the harness campaign: on its first turn a model
attempted to end the session early, the attempt was blocked by the
deny-list, and the retrieval that followed was unaffected. Sessions are
budgeted as well, with per-lane turn and call caps
(Table~\ref{tab:lanes}); the campaign ran 18 end-to-end scenarios, all of
which passed with no fabricated answer.

Three operational findings from that campaign matter to callers. A trivial
ping prompt triggered a six-skill pre-flight before any retrieval, and the
probe was killed by a 30\,s timeout; the platform's default caller timeout
is therefore 420\,s, not the few seconds a request such as ``list my
libraries'' suggests. The streaming endpoint emits bare \texttt{data:}
frames with no event names, so a client must parse data-only frames instead
of waiting for a named event. Finally, engine choice dominates cost by one
to two orders of magnitude at the same quality tier: \$0.90--1.80 per
question through the agentic surface, \$0.11--0.36 through the Claude
harness, and \$0.03--0.07 through the OMP harness (the three rows of
Table~\ref{tab:e5-cost}). Wall clock was 50--176\,s, 23--52\,s, and
34--171\,s respectively, but the third engine's engine-side time was only
15--83\,s: its extra wall clock is harness cold start of 50--90\,s per call,
not retrieval. One residual gap appeared in the same campaign: a document
ingested through the third harness had no graph node, while the same
operation through the agent surface created one.

\subsection{Implementation and runtime engineering}
\label{sec:impl-runtime}

The backend is Python 3.12 with FastAPI. ChromaDB~\cite{chroma} stores one
HNSW~\cite{hnsw} collection per knowledge base in cosine space; embeddings
are BGE-M3~\cite{bge-m3}, 1024-dimensional and normalized, encoded in
batches of 32 by a lazy-loaded singleton with an LRU encode cache. Neo4j
5.20~\cite{neo4j} holds the document graph, jieba~\cite{jieba} tokenizes the
keyword lane, MinerU~\cite{mineru} performs parsing and OCR, and MCP is
spoken over stdio and SSE.

The enforcement described in Section~\ref{sec:lanes} depends on engineering
invisible in the interfaces. The judging model runs as a resident worker
behind a machine-level file lock, with an LRU verdict cache and explicit
CUDA-cache release between workloads. A cold judge load costs 10--12\,s per
batch; the resident worker answers in 1--2\,s. The worker also removes a
failure mode met in production: when judge and bulk-embedding workloads ran
interleaved without cache cleanup, one end-to-end run degraded from 68\,s to
27 minutes under the Windows display driver model's paging behaviour. Judge
calls are serialized by one lock, because two concurrent judgments on a
single GPU raise a native allocator failure; a call is held to a 900\,s
ceiling and fails closed beyond it. MinerU
is launched as a hidden subprocess with platform-specific parent-death
cleanup---a Windows Job Object with kill-on-close, Linux
\texttt{prctl(PR\_SET\_PDEATHSIG)}, and a process group plus
\texttt{atexit} on macOS---and its stdout is redirected to a log file rather
than a pipe, because the pipe variant failed with a spurious
\texttt{Errno 22} once the reader closed. Before start, the backend probes
its own port with a bare socket bind and refuses to start if the port is
occupied, so a surviving process cannot be shadowed by a second one. Vector
writes are serialized per collection, after one production incident traced
to concurrent segment handles; the keyword index updates incrementally.
Every HTTP client is constructed with \texttt{trust\_env=False}, so a
corporate proxy cannot intercept calls to localhost.

\subsection{Configuration, deployment, and reproducibility}
\label{sec:impl-deploy}

A single \texttt{config.yml} is the source of truth for ports and
thresholds; environment variables override it. The mechanism is used in
practice, not only documented: in the harness campaign the services ran on
web 6789, backend 8771, and kb-mcp 8000, the backend port being an override
of the shipped development default. Rate limiting (a general bucket and a
smaller bucket for expensive routes), authentication tokens for the MCP,
user, and legacy clients, and the retrieval thresholds all live in the same
file. The \texttt{ragctl} command-line tool covers setup, start, status,
dependency, configuration, and update operations. The system runs
self-hosted on a single workstation: ChromaDB runs in-process, Neo4j runs
locally on bolt port 7687, and MinerU runs as a local subprocess, so no
external service is required for a deployment. The same file carries the
higher retrieval threshold (0.55) used by the experience subsystem, whose
rolling summary of operational lessons is regenerated by a rules-based pass
every fifth ingestion. That loop is implemented but not evaluated:
Section~\ref{sec:eval}'s experience evidence is one scenario, and scheduled
distillation at scale, approval, decay and application feedback remain
unexercised (Section~\ref{sec:e3}). Source code and all experiment scripts are public at
\url{https://github.com/kingdol666/rag-knowledge}.

\section{Evaluation}
\label{sec:eval}

We evaluate the deployed system on three corpora that differ in origin, language and scale, using three metric families, and we organise the evidence into eight groups. E1 and E2, reported here, cover the stages that precede a question: acquisition and parsing, and organisation, consistency and management operations. E3--E8, reported in Sections~\ref{sec:e3}--\ref{sec:e8}, cover retrieval quality and the role of the verdict, the operational-knowledge store, end-to-end agent behaviour, verifier failure, non-functional behaviour, and the component ablations. Numbers are reported against the corpus that produced them, and results from different corpora are never merged. The complete campaign inventory---corpus, sample size, metric and artifact path for every system-level and retrieval test summarised here---is provided as supplementary material, and per-question tool-layer results are given in Appendix~\ref{app:perq}.

\subsection{Experimental setup}
\label{sec:eval-setup}

\textbf{Corpora.} (i)~\emph{Deployed industrial corpus}: five knowledge bases in a live deployment hold 42 documents, one of them containing twenty near-duplicate equipment-telemetry reports that share prefixes, vocabulary and structure, a deliberate stressor for near-duplicate detection and confuser density~\cite{manku-near-dup-2007}. Sixteen gold questions and four out-of-corpus probes were run against it. (ii)~\emph{arXiv corpus}: 100 submitted papers spanning about 25 domains; at evaluation time the state file records 91 paper names and 79 papers are retrievable. Content routing placed them in five domain libraries holding 322 stored parts, and ten gold questions (BQ01--BQ10) target books referenced by these papers, keyed by arXiv identifier. The live environment held 19 libraries. (iii)~\emph{Cross-lingual ingestion corpus}: 400 Wikipedia pages per language for English, Chinese and Japanese, with 40 extractive cloze queries per language whose answers occur verbatim in the source page. It is used for ingestion only; we report no retrieval-quality numbers for it.

\textbf{Questions and gold annotations.} The sixteen industrial questions (Table~\ref{tab:questions}) were written by the authors rather than sampled from logs, a circularity we state rather than hide; each identifies one or two gold documents by a unique filename fragment and one to three fact markers. Q16 names the event \texttt{heatdecay092}, a token that also occurs in its gold document's filename, so a retriever ranking on filenames alone could match it without reading content; we keep the question and report the weakness. On the arXiv corpus a question passes only if the gold identifier is traceable in the answer or the tool trajectory and at least 60\% of the gold keywords occur at word boundaries, evaluated deterministically with no language model in the loop.



\textbf{Metric families.} Retrieval results use Hit@5, Recall@5 and MRR@5, following passage-ranking evaluation conventions~\cite{craswell-msmarco-2021}. A graded measure, nDCG@5~\cite{jarvelin-cumulated-gain-2002}, was computed only inside one exploratory arm of the cross-lingual benchmark that we do not report; no nDCG@5 value appears here. Grounding and key-fact verification use deterministic string tests: whether the answer cites the gold document, whether at least 60\% of gold key facts appear at word boundaries, and whether part, section and table citation markers are present. These checks are mechanical, not model-based faithfulness estimates such as RAGAS~\cite{es-ragas-2024} or ARES~\cite{saad-falcon-ares-2024}. The third family covers wall-clock time, agent turns, tool calls, monetary cost, and the correctness of refused answers.

\textbf{Controls.} Deterministic lanes were re-run and compared field by field with latency excluded, so a re-run could disagree only on content and ordering; model-driven lanes were repeated with the same model version and are reported as mean with the observed spread. A fixed retry policy applied: an empty zero-turn response or a transient gateway timeout was retried once, which fired on two runs and succeeded both times. Before measurement we repaired two corpus-integrity faults through the platform's own operations: one document missing from its vector collection was re-indexed, and one library whose collection had been silently lost was rebuilt. All measurements ran on one workstation whose RTX~5060 (8\,GB) serves the embedding model and the judge worker, with every service local under Windows.

\textbf{What is deterministic.} The tool-layer retrieval lanes, the grading rule, the audits, duplicate adjudication and the cross-lingual ingestion pipeline are deterministic: repeating them reproduces the same values, and cross-lingual round~2 was bit-identical to round~1. The agent lanes, and anything derived from generated text, are not, and are reported with their repetition spread.

\subsection{E1: Acquisition and parsing}
\label{sec:e1}

One hundred submitted papers were parsed and routed by content into five domain libraries, producing 322 stored parts after the 30{,}000-character split. The funnel is reported with the audit rather than beside it: 91 papers reached the ingest state file and 79 were retrievable at evaluation time. Table~\ref{tab:acq-pipeline} lists the acquisition audit: nine per-paper ingest checks passed 100/100, no true duplicate was found against 35 near-duplicate pairs, 165 document descriptions were upgraded, the content and storage audits passed 322/322 and 33/33, all 15 retrieval regression probes passed, and the dense replica index holds 4{,}528 chunks. Values are asserted by the run scripts from saved artifacts rather than hand-counted; the content audit samples the first 1{,}000 characters of each part and the storage audit covers 33 storage units.

\begin{table}[t]
\centering
\small
\setlength{\tabcolsep}{4pt}
\caption{Acquisition and parsing audit over the 100 submitted arXiv papers.
The first four rows are the submission funnel: every paper was parsed and
routed by content, but only 91 reached the ingest state file and 79 were
retrievable at evaluation time, a gap we cannot attribute from the artifacts.
The per-paper check row counts parse tasks; the funnel rows count papers the
state file records, and the two ledgers differ by the nine names it omits.
Values are asserted by the run scripts from saved artifacts, not hand-counted.
The content audit samples the first 1{,}000 characters of each part. Duplicates
are compared by fingerprint at a 0.90 threshold~\cite{manku-near-dup-2007}.}
\label{tab:acq-pipeline}
\begin{tabular}{@{}lr@{}}
\toprule
Check & Result \\
\midrule
Papers submitted to the pipeline & 100 \\
Papers parsed and routed by content & 100 \\
Papers recorded in the ingest state file & 91 \\
Papers retrievable at evaluation time & 79 \\
Content-routed libraries & 5 \\
Stored parts after 30{,}000-character split & 322 \\
Per-paper ingest checks (nine checks, per parse task) & 100/100 pass \\
True duplicates (35 near-duplicate pairs) & 0 \\
Descriptions upgraded & 165/165 \\
Content audit & 322/322 \\
Storage audit & 33/33 \\
Retrieval regression probes & 15/15 \\
Dense replica index (800-character chunks) & 4{,}528 chunks \\
\bottomrule
\end{tabular}
\end{table}

\textbf{Accounting gap.} This campaign does not support the claim that all 100 papers became usable, and we do not make it: 100 papers were submitted, the state file \texttt{exp10\allowbreak\_ingest\allowbreak\_state.json} records 91 names, and 79 were retrievable at evaluation time, so 21 papers were never retrievable. The reported failure modes are formula-dense scanned papers, which the parser handled poorly, and one crash of the tool layer during batch parsing that was restarted from a checkpoint. No artifact records a count of attempted against failed papers, and the state file has no per-paper parse status, so we cannot say whether the nine papers missing between submission and the recorded list failed or were never recorded. The absent counter is an observability gap: reconciliation is possible only against the artifacts that survive~\cite{missier-prov-2013}.

\textbf{Cross-lingual ingestion.} The benchmark ran twice over 400 pages per language, with both rounds retained under \path{benchmark-web/benchmark/results/benchmark/}. The ingest success rate was 1.0 in both rounds for all three languages. Membership accuracy, whether a stored page landed in the intended library, was 0.875 for English, 1.0 for Chinese and 1.0 for Japanese (400 pages per language). Storage completeness, computed on a sample of 40 pages per language with the shards of large documents summed, was 0.975, 0.9559 and 0.9402 respectively. Self-retrieval Hit@1, again at 40 pages per language, was 0.975, 1.0 and 1.0. Round~2 was bit-identical to round~1. The English figure of 0.875 is the weakest cross-lingual number we report, and it is an organisation defect rather than a parsing defect: the page was stored, but in the wrong library roughly one time in eight.

\textbf{Smaller campaigns.} A 50-paper, 26-domain campaign parsed and verified 50/50 papers, took 1{,}115.5\,s at 16.3--40.4\,s per paper, produced 202 tags, batch-indexed 50 documents and reported \texttt{all\_probes\_ok}; in its routing run 43 of the 50 documents carried \texttt{tag\_ok:false} while \texttt{verify\_ok} passed for all 50, so the tag-quality check failed often without blocking ingestion. Its 30{,}000-character tier turned 50 papers into 165 documents (nine whole documents and 156 parts), passed the storage check 50/50 and the description-and-tag check 165/165 with zero failures, stored 3{,}798{,}446 characters and took 2{,}672\,s to parse. Two caveats attach: one stage re-ran because of inbox ordering, and more than 800 extracted images were not uploaded. Descriptions and tags are checked against document content before storage~\cite{park-metadata-quality-2009}.

\subsection{E2: Organization, consistency, and management operations}
\label{sec:e2}

We define organisational correctness as three-source consistency: the files on disk, the tree index and the per-library metadata must agree about which documents exist, where they live and how much indexed content each has. Under that definition we found a real defect in a disposable test library. After a document was renamed, the vector index retained the chunks of the old name: the recorded chunk count was 18 where 9 were expected, a duplicate search matched the document against its own former path at 0.9738 similarity, and a vector probe returned every chunk twice with identical scores. Neither the automatic re-indexing that follows a rename nor a forced batch index removed the stale chunks; only an explicit force re-index repaired the document, after which the recorded count was 9. The failure is invisible from the metadata alone.

\textbf{Reorganisation campaign.} The full reorganisation workflow then ran over the 322-document arXiv corpus: 322/322 metadata checks and 165/165 description checks passed, the tag registry was de-duplicated, all 322 documents were batch-indexed, 15/15 retrieval regression probes passed with zero degradation, and 35 near-duplicate candidate groups (similarities 0.90--0.975) were adjudicated as zero true duplicates, every group being a set of adjacent parts of the same paper; no document was deleted (Table~\ref{tab:org-campaign}). Two honesty notes apply: the three-source spot check sampled 33 of the 322 documents, so it is a sample rather than a full pass, and sub-library splitting and hierarchy creation were analysed but deliberately not executed, pending user confirmation, so the campaign is not evidence for hierarchical restructuring.

\begin{table}[t]
\centering
\small
\setlength{\tabcolsep}{2pt}
\caption{Reorganisation campaign over the 322-document arXiv corpus.
The three-source spot check samples 33 documents; the description, tag and
duplicate checks cover the full corpus.}
\label{tab:org-campaign}
\begin{tabular}{@{}lrr@{}}
\toprule
Check & Scale & Result \\
\midrule
Metadata consistency (three sources) & 322 documents & 322/322 pass \\
Description check & 165 documents & 165/165 pass \\
Tag de-duplication & 4 groups & 4/4 pass \\
Batch index & 322 documents & 322 indexed \\
Duplicate adjudication & 35 groups & 0 true duplicates \\
Three-source spot check & 33 of 322 documents & 33/33 pass \\
Retrieval regression probes & 15 probes & 15/15, no degradation \\
\bottomrule
\end{tabular}
\end{table}

\textbf{Scale.} Growing the corpus from 50 to 100 papers (163 to 322 documents) did not cost retrieval quality: dense top-1 hits went from 48/50 to 49/50 and MRR from 0.98 to 1.00, and reorganisation itself produced a 2 percentage point improvement. The judgement field recorded with that comparison reads ``held (within $\pm$5\,pp)''; we report the field rather than restating the gain as a retrieval improvement.

\textbf{Lifecycle operations.} Table~\ref{tab:org-ops} lists the operations exercised. Create, move, delete and forced re-index behaved as specified: the move left the source library at zero documents and the target at one, and batch deletion removed both disposable test libraries. Two failure modes are worth reporting. The task-status tool reports a task as \texttt{running} after it has finished: the record carries a \texttt{finished\_at} timestamp while the status field still reads \texttt{running}, and the same symptom appeared in the reorganisation graph build, where node reconciliation replaced task polling. The update check timed out after 30\,s (\texttt{ragctl version -{}-json timed out after 30s}, success false) while every other probe in the batch passed; the project-status tool also silently ignores its version flag. Merge and split were never exercised under controlled conditions and we claim nothing about them. These 16 skill scenarios ran against two disposable test libraries in a standalone skill environment; the ingest tools were exercised in E1, not here.



\textbf{Tags and graph.} The live tag registry held 316 tags, and the tag-quality check passed 165 documents with zero failures in 47.6\,s. Orphan tags do accumulate, among them test tags and section titles extracted from headings; the registry offers a dry-run detector and a cleanup operation, and in one case tags from a test library leaked into the production registry, the failure mode that vocabulary-based organisation is prone to~\cite{golder-tagging-2006}. The document graph is used for navigation and candidate expansion, never as evidence. Writing one document's four tag edges moved the edge count from 4{,}675 to 4{,}679, against 658 nodes, 333 documents and 316 tags; graph statistics for the 165-document corpus were 335 nodes and 1{,}671 edges.

%% n8b_eval.tex -- evaluation, middle part: E3 (operational knowledge and
%% automated distillation) and E4 (retrieval quality and the role of the
%% verdict). Assembled 2026-10-12. Labels owned here: sec:e3, sec:e4,
%% tab:retrieval, tab:e4-modes.
%% Numbers: restructure-plan-2026-10-12.md (\S C.1 items 8.4/8.5, \S D-3, \S H),
%% functional-test-evidence-2026-10-12.md, v2/s7_OLDresults.tex, and the
%% artifacts named inline (paper-kbs/experiments/results/, benchmark-suite/
%% results/, review/fulltest-20261012, review-2026-10-08.md).
\subsection{E3: Operational knowledge and automated distillation}
\label{sec:e3}

The experience store holds operational lessons as first-class retrievable
documents: an entry is addressed, indexed and searched through the same
interfaces as an ordinary document, and it carries a credibility tier that
decides whether it may be shown at all
(Section~\ref{sec:org-experience}, Table~\ref{tab:exp-tiers}). The design
claim is that a lesson about practice should be governed like any other
content object. The evidence collected for that claim is case evidence and we
report it as such.

Three measurements exist. (i)~One extraction campaign over a single small
knowledge base produced thirteen candidate drafts
(\path{benchmark-suite/results/experience_smoke.json}:
\texttt{extract\_candidates}~13, \texttt{drafts\_created}~13). (ii)~One
end-to-end scenario was sedimented into one experience
(\texttt{exp-d6f4d47ae3d6}) and a later global experience search returned that
entry first, with a local similarity of 0.711
(\path{review/fulltest-20261012}, scenario A6). (iii)~The draft-approval
path was never exercised: the draft pool was empty when it was queried
(\texttt{experience\_drafts\_list} returned \texttt{count}~0 with no drafts),
so \path{experience_draft_read} and \path{experience_draft_approve}
could not run at all (\path{benchmark-suite/results/SKILL-E2E-REPORT.md},
lines~27 and~103). The heuristic extraction path was exercised only in
dry-run mode, where it returned a raw document dump that the quality gate
would reject (line~100).

One scenario and thirteen candidates cannot support an effectiveness claim,
and we do not make one. The larger part of the automated-distillation
pipeline is unverified: scheduled distillation, review and approval at scale,
application feedback, and credibility decay were not exercised end to end.
What the measurements show is narrower and sufficient for the design
argument: an operational lesson can travel the document path --- created,
indexed, returned first by search, with its tier attached --- and nothing more.

\subsection{E4: Retrieval quality and the role of the verdict}
\label{sec:e4}
\label{sec:eval-retrieval}

Nine configurations, defined here so the table can be read on its own. Seven
were measured through their own tool-layer entry points: \emph{BM25} through the
keyword search tool; \emph{Dense} through \texttt{kb\_search\_vector} with
BGE-M3 \cite{bge-m3}; \emph{RRF} as reciprocal-rank fusion at $k=60$ of those
two rankings \cite{rrf-cormack}; \emph{Graph} through
\texttt{kb\_graph\_search}, a node-name and path match over the document graph;
\emph{Two-stage} through \texttt{kb\_search\_two\_stage}
(Section~\ref{sec:twostage}); and the hybrid lane through
\texttt{kb\_hybrid\_search} with the gate switched off
(\emph{Fusion-only}) and on (\emph{QDCVR-Hybrid}). That lane is implemented in
the shipped tool layer as a server-side composition of the vector, catalog and
judging tools named above; \texttt{kb\_hybrid\_search} was an earlier single-call
wrapper that is no longer registered. Two further rows are derived
offline by replaying the judged hybrid run's own logged twelve-candidate pools
and verdict scores: \emph{Judge-as-filter} keeps the survivors in their fusion
order, and \emph{RRF-fuse} fuses the fusion ranking with the judge-score
ranking at $k=60$. The last three configurations see the same candidates and
the same scores; only the verdict's role differs, which is what makes the
comparison an attribution rather than a benchmark.

\begin{table}[!htbp]
\centering
\small
\setlength{\tabcolsep}{3pt}
\caption{Tool-layer retrieval on the 42-document deployed industrial corpus
(mean over 16 gold-annotated questions; top-5 document lists; latency per
call, warmed). The paraphrase column re-asks those sixteen questions in
different words (\protect\path{paper-kbs/experiments/results/exp6_paraphrase.json});
block~(b) repeats the Hit@5 measurement on a second corpus, the 79 retrievable
papers of the 100-paper arXiv ingestion, over ten questions.}
\label{tab:retrieval}
\resizebox{\textwidth}{!}{%
\begin{tabular}{lcccccc}
\toprule
Method & Hit@5 & Recall@5 & MRR@5 & Paraphrase Hit@5 & Latency (mean) & Latency (median) \\
\midrule
BM25                  & 1.00 & 1.00 & 0.958 & 0.88 & 0.9\,ms    & 1\,ms \\
Dense (BGE-M3)        & 0.94 & 0.94 & 0.856 & 0.94 & 18\,ms     & 17\,ms \\
RRF (BM25$\|$Dense)   & 1.00 & 1.00 & 0.927 & 0.88 & 11\,ms     & 9\,ms \\
Graph (Neo4j)         & 0.56 & 0.53 & 0.562 & 0.56 & 68\,ms     & 68\,ms \\
Two-stage             & 1.00 & 0.97 & 0.927 & 0.94 & 77\,ms     & 40\,ms \\
Fusion-only           & 0.94 & 0.94 & 0.856 & 0.94 & 62\,ms     & 66\,ms \\
QDCVR-Hybrid          & 0.44 & 0.41 & 0.226 & 0.50 & 542\,ms    & 582\,ms \\
\midrule
\multicolumn{7}{l}{\emph{Derived (offline, same logged hybrid pools):}}\\
\quad Judge-as-filter$^{\dagger}$ & 0.94 & 0.94 & 0.856 & 0.94 & \multicolumn{2}{l}{$^{\dagger}$542\,ms / 582\,ms} \\
\quad RRF-fuse & 0.88 & 0.88 & 0.606 & 0.88 & \multicolumn{2}{l}{(same single call)} \\
\bottomrule
\end{tabular}}
\par\smallskip
\raggedright\small
$^{\dagger}$\,Derived rows replay the QDCVR-Hybrid run's own logged
12-candidate pools and judge scores, so they reuse that call's latency; the
judge executes once either way.

\par\smallskip
\resizebox{\textwidth}{!}{%
\begin{tabular}{@{}lcccccccc@{}}
\toprule
\multicolumn{9}{@{}l}{\textbf{(b)} Second corpus: 79 retrievable arXiv papers of the 100 submitted, ten questions}\\
\midrule
Configuration & BM25 & Dense & RRF & Graph & Two-stage & Fusion-only & QDCVR-Hybrid & Judge-as-filter \\
Hit@5 & 0.90 & 0.90 & 0.90 & 0.20 & 0.90 & 0.90 & 0.90 & 0.90 \\
\bottomrule
\end{tabular}}
\end{table}

Table~\ref{tab:retrieval} compares nine configurations on the sixteen
gold-annotated questions of the deployed industrial corpus (42 documents;
\path{paper-kbs/experiments/results/exp1_retrieval.json} and
\texttt{exp5\_judge\_filter.json}). The tool-layer arms retrieve and never
answer: their latencies are retrieval latencies, and, because no generation
follows, they carry no token cost. With sixteen questions the Wilson 95\%
intervals on the headline proportions are wide (judged hybrid $[0.23, 0.67]$;
judge-as-filter $[0.72, 0.99]$; BM25 and RRF $[0.81, 1.00]$), so the
decomposition below, not the point estimates, is the claim. A separate
six-method retrieval-mode ablation exists in the released logs but covers two
questions; it is a smoke check at $n = 2$ and no claim rests on it. The
2026-10-12 re-run of that ablation
(\path{benchmark-suite/results/runs/run-20261012T062847Z-5f9faac/RETRIEVAL-COMPARE.md})
returns the gold in all twelve method--question cells, with mode~B at
25.0/10.2\,s and mode~C at 37.3/33.7\,s where the 2026-09-27 comparison has
them at 78.5/29.8\,s and 235.1/123.4\,s
(\path{benchmark-suite/results/runs/run-20260927T094215Z-47cc948/RETRIEVAL-COMPARE.md}),
and an earlier 2026-09-26 snapshot --- eleven libraries, before a catalog
change --- recorded $\sim$342/234\,s and $\sim$1152/1818\,s for the same two
modes
(\path{benchmark-suite/results/runs/retmodes-20260927T062953Z-47cc948/THREE-MODE-REPORT.md}),
though its mode~C arm is not contract-matched, keeping 173 surviving documents
against 36 now.

\textbf{On this corpus, lexical retrieval is strong.} BM25 \cite{bm25}
reaches Hit@5 1.00 (MRR 0.958) at about a millisecond per call; reciprocal
rank fusion \cite{rrf-cormack} matches it in Hit@5 and follows in MRR (0.927);
the two-stage BM25+graph-then-dense pipeline reaches Hit@5 1.00 with
Recall@5 0.97, one gold of a two-gold question falling outside the top-5.
Dense retrieval with BGE-M3 \cite{bge-m3} misses exactly one question (Q16),
where four near-duplicate sibling reports from one incident family crowd the
top-5 --- a confuser-density failure, not a vocabulary one. Graph search over
the document graph \cite{graphrag}, limited to node-name and tag matching, is
the weakest standalone signal (Hit@5 0.56, MRR 0.562). Nothing here is a
retrieval advance: the corpus is curated, the questions are written in its
vocabulary, three configurations tie at Hit@5 1.00, and the one that uses only
words leads on MRR.

\textbf{The central result is the verdict's role, not its existence.} The same
judged scores enter three configurations: not used (fusion-only, Hit@5 0.94),
used to re-rank candidates (QDCVR-Hybrid, 0.44), and used as a filter that
keeps survivors in their fusion order (judge-as-filter, 0.94). On the eight
questions where re-ranking and filtering differ (Q2, Q4, Q5, Q6, Q8, Q9, Q11,
Q12), filtering wins all eight and loses none, with eight further ties on
Hit@5, and it restores MRR to 0.856 (exact binomial sign test, two-sided
$p = 0.008$); the only ranking movement is Q7, where the gold slips from third
to fifth inside the top-5. The verdicts were not wrong about which documents
are evidence --- they were wrong as an ordering. This matches the
judge-miscalibration pattern reported for model-generated preferences
\cite{zheng-llmjudge,panickssery-selfpref}: an individually plausible verdict
can be ordinally unreliable, and the fix is the placement of the verdict, not
its existence or its strength. Its price is measurable: judging
twelve candidates adds well under a second with the resident worker, and the
hybrid call's latency surcharge over fusion-only is the observed difference
between 542\,ms and 62\,ms per call. Two caveats travel with the test. It is
post-hoc and descriptive --- one of nine configurations evaluated on the same
sixteen questions --- and the gate is a single locally fine-tuned checkpoint
(Section~\ref{sec:gate}), so every verifier finding in this paper inherits
that choice.

\textbf{A second corpus reproduces the baselines, not the collapse.} Ten
questions over the 79 retrievable papers of the 100-paper arXiv ingestion
(the campaign record logs 8{,}031 chunks; block~(b) of
Table~\ref{tab:retrieval}; per-question results in
\path{paper-kbs/experiments/results/exp10_papers_toollayer.json}) are answered
at Hit@5 0.90 by seven of the eight configurations, graph search excepted at
0.20. Question BQ03 is missed by all eight configurations, and apart from
graph search no other configuration-level difference is worth noting: its gold
paper is one of the 21 papers that never became
retrievable in that campaign, so the single shared miss is an acquisition-side
loss, not a retrieval failure. We therefore report both denominators --- 0.90 over
the ten questions as asked, 1.00 over the nine answerable ones --- and say
which is which. The judged collapse of the industrial corpus does not
reproduce on this corpus, where the same-family confuser cluster that drives
it is absent; the graph configuration's 0.20 tracks absent relation tagging,
since the bulk import bypassed the tagging path.

\textbf{Thresholds and margins.} All 191 judged candidates from the sixteen
hybrid runs passed the 0.5 acceptance threshold, the lowest logged score being
0.589, so at the threshold layer the gate filtered nothing on this corpus
(\texttt{exp1\_retrieval.json}). The margin sweep that follows replays these
logged candidate pairs rather than re-retrieving them: with the top-5 floor
unchanged, a margin of 0.15 keeps the gold in 16 of the 17 in-pool gold
instances and 0.20 keeps it in all 17
(\texttt{exp6\_loo\_margin.json}). A leave-one-out procedure that
selects the margin per question from the other fifteen reaches a gold survival
of 0.9375, its single miss being Q11, which the shipped 0.15 margin also
loses, while a fixed 0.20 floor reaches 1.00
(\texttt{exp6\_loo\_margin.json}). The shipped floor is thus an engineering
constant, not a fitted parameter, and the calibration that would justify it
has not been run.

\textbf{Rephrasing shows where the fragility sits.} On a grid of sixteen
paraphrase questions (\texttt{exp6\_paraphrase.json}) the judged configuration
falls to Hit@5 0.5 (MRR 0.189), while judge-as-filter and fusion-only hold at
0.9375 (MRR 0.833) and dense retrieval stays at 0.9375; BM25 falls from 1.00
to 0.875 (MRR 0.958 to 0.755). Both the lexical baseline and the judged
re-ranking lose ground when the wording changes, and the configuration that
refuses to let the verdict order the list does not. The gate's weakness is in
the ordering it induces, not in the verdicts it returns.

\textbf{End to end, the modes differ in checkability.} Ten fixed questions ran
through three modes that share one harness, one model and one neutral prompt
template and differ only in retrieval affordance: the platform's full read
surface, a bare agent with file tools over the exported papers, and dense
retrieval over an 800-character replica index of the corpus (4{,}528 chunks).
Table~\ref{tab:e4-modes} reports the outcome. All three modes produce answers,
and the bare agent grounds as many of them as the platform
(10/10 vs.\ 10/10), but its answers name files where the platform names
checkable addresses, and dense retrieval cites an address once in ten
questions. The modes differ mainly in checkability rather than in whether an
answer was produced. Three caveats belong to this comparison: both dense
grounding misses were tool-availability aborts; grading is conservative, and
word-boundary matching zeroes the bare agent's plural ``cathodes'' on BQ08; and
two platform cells re-ran once after a harness timeout, correctness calls are
the authors' trace readings, and cache state was uncontrolled.

\begin{table}[!htbp]
\centering
\small
\setlength{\tabcolsep}{4pt}
\caption{Three-mode end-to-end comparison on ten fixed questions over the
100-paper arXiv corpus. \emph{Grounded}: strict arXiv-id citation of the
designated paper. \emph{Key-fact}: at least 60\% of the fixed gold key facts
matched at word boundaries. \emph{Cited}: the answer names a part, section or
table. Wall clock and cost are per question.}
\label{tab:e4-modes}
\begin{tabular}{lccccc}
\toprule
Mode & Grounded & Key-fact & Cited & Wall clock & Cost \\
\midrule
Platform (read surface)   & 10/10 & 10/10 & 9/10 & 116\,s & \$0.55 \\
Bare agent (file tools)   & 10/10 & 8/10  & 6/10 & 48\,s  & \$0.63 \\
Dense retrieval (replica) & 8/10  & 6/10  & 1/10 & 41\,s  & \$0.28 \\
\bottomrule
\end{tabular}
\end{table}

E4 establishes one thing. Content verification placed as a filter rather than
as a ranker is the difference between a citation and a claim: every returned
document is judged, and the ordering the reader sees is not the verdict's.
Verification before answering is not itself new \cite{selfrag,crag}; what is
measured here is where the verdict belongs in the pipeline and what it costs
there. E4 does not establish a retrieval advance --- BM25 reaches Hit@5 1.00
on this corpus (Figure~\ref{fig:retrieval}) --- so recall is not the contribution, and no margin, threshold
or paraphrase variant tested repairs the ordering role of the verdicts. That
is the question the failure analysis of the verifier takes up.
Figure~\ref{fig:marginvariants} plots both variants on the same logged
candidate pairs, so the margin rule and the paraphrase collapse can be read
against the aggregate numbers above.

\begin{figure}[t]
\centering
\includegraphics[width=\textwidth]{fig3_retrieval}
\caption{Tool-layer retrieval over the 42-document production corpus
(sixteen gold questions, top-5 document lists). Lexical retrieval is
strong on this corpus: BM25 reaches Hit@5 1.00. The judged hybrid loses
more than half of its top-5 hits (0.44); placing the same verdicts in a
filter role recovers the fusion result (0.94) because the verdicts were
not wrong about which documents are evidence, only about their order.}
\label{fig:retrieval}
\end{figure}

\begin{figure}[t]
\centering
\includegraphics[width=\textwidth]{fig6_margin_variants}
\caption{Margin and paraphrase variants over the same 191 logged judged
candidate pairs. (a) The relative cut replays the logged pairs with the
top-5 floor unchanged: the shipped 0.15 margin loses one of seventeen
in-pool gold instances, a 0.20 floor loses none, and a leave-one-out
selection over the other fifteen questions lands at 0.9375. (b) Under
sixteen paraphrases the judged re-ranking collapses to Hit@5 0.50
(MRR 0.189) while the filter repair and dense retrieval hold at 0.9375
and BM25 falls to 0.875 (MRR 0.755).}
\label{fig:marginvariants}
\end{figure}

% n8c_eval.tex — evaluation 8.6-8.8: E5 agent behaviour/abstention/cost,
% E6 verifier failure, E7 non-functional behaviour.
% Every figure is attributed to its artifact; corpora are never mixed.

\subsection{E5: End-to-end agent behaviour, abstention, and cost}
\label{sec:e5}
\label{sec:agentic}\label{sec:abstention}\label{sec:composition}\label{sec:eval-agent}
%% The four labels above are compatibility aliases: n2_related.tex:110,122,139
%% and n6_retrieval.tex:38,53-54 were written against the pre-restructure label
%% set. They all resolve to this subsection (agent behaviour, abstention,
%% latency composition).

\textbf{The agentic grid.}
Four arms --- the Search, Librarian and Hybrid lanes plus a bare model with
no retrieval tools --- answered four questions (Q2, Q4, Q6, Q7) twice each
through the public chat API, one engine (\texttt{omp} over GLM-5.3-Flash),
20-turn budget. Citation outcomes agree on all 16 of 16 arm--question cells:
every grounded cell named the gold document (12/12); the bare arm cited a
source in 0 of 4 cells in either round. The de-contaminated fact-marker
check (delete answer lines naming a gold document from the first 600
characters, re-test the remainder) is 12/12 in round 1 and 11/12 in round 2;
the exception, Hybrid-Q2, keeps its markers only inside lines that also cite
the gold, an artifact of the deletion rule. Six of the twelve grounded cells
answer from libraries of one or two documents, where hitting the gold is a
low bar; the substantive evidence is the six cells over the 31-document
process library. Where the bare arm is wrong, it is
wrong in the way that motivates retrieval: on the lamination question it
gave a plausible textbook diagnosis (interfacial adhesion failure, moisture
ingress) missing the corpus's recorded root cause, thermocouple drift
leaving the encapsulant under-cured.

\textbf{Abstention is a property of the reading-and-answering layer, not of
a threshold.}
Three out-of-corpus questions produced three explicit not-found reports
(47.9, 73.3 and 113.5\,s wall clock; turns were not recorded for this probe
family; \path{paper-kbs/experiments/results/exp8b_absent_agent.json}), each
enumerating the scope it searched (five libraries, 42
documents). The retrieval tools themselves never abstain: six graded probes
at the tool layer
(\path{paper-kbs/experiments/results/exp8_graded_probes.json};
$n = 6$; three partially covered, three absent) retained 10
to 12 of their 12 candidates at every tier, with top scores up to 0.983 on
absent probes and 0.965 on partial ones; the threshold of 0.5 never fired as
a refusal. Honesty is enforced where the answer is written, in the lane's
read-then-answer contract \cite{selfrag,feng-abstain,yin-knowwhat}, not by a
similarity cut-off.

\begin{figure}[t]
\centering
\includegraphics[width=0.62\textwidth]{fig7_ood}
\caption{Out-of-corpus probes at the tool layer
(\protect\path{paper-kbs/experiments/results/exp3_abstention.json}): for every probe
the judged hybrid keeps 11, 11, 9 and 11 of the 42--43 candidates its lanes
merged and dense
retrieval passes chunks above the 0.35 threshold (top vector scores
0.397--0.596, above it for 4, 8, 10 and 7 of the graded pool), so the
threshold alone never abstains and the 0.5 refusal threshold never fires. At
the agent layer, three absent-tier probes --- a separate family
(\protect\path{paper-kbs/experiments/results/exp8b_absent_agent.json}) --- each
returned an explicit not-found report with zero fabricated citations.}
\label{fig:ood}
\end{figure}

\textbf{Where the time goes.}
In the twelve grounded cells, machine-side retrieval and verification
stayed sub-second (BM25 0.9\,ms mean, dense 18\,ms, two-stage 77\,ms,
judged hybrid 542\,ms mean) against an end-to-end wall clock of 36--91\,s.
The gap is model turns: 11.7, 12.8 and 15.3\,s per turn for the Search,
Librarian and Hybrid arms --- the turn-dominated profile reported for
agentic retrieval \cite{singh-agentic-rag}, and the reason orchestration
runs server-side with a turn cap. Results repeat; wall clock does not:
Search-Q2 took 70\,s over 6 turns in round 1 and 165\,s over 11 turns in
round 2, Librarian-Q4 84\,s over 6 turns and 92\,s over 11
(Figure~\ref{fig:e2e}).

\textbf{Cost.}
The same task set on three engines separates by one to two orders of
magnitude at one quality tier (Table~\ref{tab:e5-cost}): all 18 scenarios
pass with every answer gold-hit and no fabricated citation. The cheapest
engine pays a 50--90\,s harness cold start per call, so its wall clock is
about twice its engine time; the Claude harness's wall clock is
approximately its engine time.

\begin{table}[!htbp]
\centering
\small
\setlength{\tabcolsep}{4pt}
\caption{Three engines on one task set (18 runs) and the two model-backed
delivery channels of the 2026-10-12 three-channel campaign (9 runs each;
\protect\path{review/retrieval-e2e-20261012/channel1_summary.json} and
\protect\path{benchmark-suite/results/runs/chatmodes-20261012T061102Z-5f9faac/CHAT-MODES-COMPARE.md}).
Cost is per run in USD as metered by the harness; the non-streaming chat
channel does not meter cost.}
\label{tab:e5-cost}
\begin{tabular}{p{4.4cm}rrrp{2.4cm}}
\toprule
Engine or channel (surface) & Runs & Wall (s) & Turns & Cost per run \\
\midrule
Agent surface (\texttt{/api/kb/agent/chat}) & 7 & 50.5--175.9 & 2--9 & \$0.9--1.8 \\
Claude harness (\texttt{engine=claude}) & 6 & 23.2--52.4 & 5--8 & \$0.11--0.36 \\
OMP harness (\texttt{engine=omp}) & 5 & 33.5--171.4 & 2--9 & \$0.03--0.07 \\
\midrule
\multicolumn{5}{l}{\emph{Delivery channels, the same three modes (9 runs each):}}\\
Channel 1 (streaming harness, \texttt{engine=claude}; 8/9 gates) & 9 & 29.9--353.5 & 5--8 & \$0.12--0.25 \\
Channel 2 (non-streaming chat API; 8/9 gates) & 9 & 25.7--100.0 & 4--11 & -- \\
\bottomrule
\end{tabular}
\end{table}

\textbf{Three delivery channels, one shared failure.}
A timestamped follow-up campaign (2026-10-12) ran the same three modes
through three delivery channels
(\path{review/retrieval-e2e-20261012/REPORT.md}). Channels 1 and 2 --- the
streaming harness (\texttt{engine=claude}) and the non-streaming chat API ---
each pass 8 of 9 gates, and their single failure is the same cell, B$\times$cm1,
a false refusal reproduced independently in both. In the streaming channel
that cell spent 300.0\,s inside a single \texttt{kb\_search\_vector} call
(353.5\,s wall), so a stall was reported as absence rather than retrieved
around. Channel 3 removes the model layer: pure retrieval returns the gold in
all 12 method--question cells, the same three modes run in 10.2--37.3\,s ---
an order of magnitude below the slowest streaming arm --- and the lexical and
vector baselines finish in 0.0--0.1\,s at zero token cost
(\path{benchmark-suite/results/runs/run-20261012T062847Z-5f9faac/RETRIEVAL-COMPARE.md}).
We report this as a dated observation on this corpus, and the gate tally uses
the campaign's corrected refusal-priority reading, under which an answer that
is substantively a refusal to an answerable question fails.

\subsection{E6: When the verifier fails}
\label{sec:e6}
\label{sec:inflation}\label{sec:eval-verify}   %% compatibility aliases (n2_related.tex:110)

\textbf{Same-domain score inflation.}
The 191 candidates judged across the sixteen hybrid tool runs --- twelve
per run except Q14, whose merge produced eleven --- carry no usable
separation between gold and non-gold. Every candidate scored at least
0.589, far above the 0.5 acceptance threshold, so on this corpus the
threshold rejects nothing. The 17 in-pool gold instances (Q16's gold never
entered its pool) have a median of 0.88 and the 174 non-gold candidates
0.86, a difference a two-sided Mann--Whitney test cannot distinguish from
chance ($p = 0.26$; Figure~\ref{fig:inflation}). Only five of the 17 golds rank in the top three of
the judged ordering and nine rank eighth or lower. Re-ranking with the
verdict therefore loses gold: the judged ordering reaches Hit@5 0.4375
(7/16; 0.44 in Table~\ref{tab:retrieval}), where the verdict as a filter,
or within rank fusion, holds 0.94 (Section~\ref{sec:e4}).

\textbf{The verdict is a model \cite{laya}, and a second model fails the
same way.}
The same 191 logged (query, candidate) pairs --- a replay, not a
re-retrieval --- were replayed through a generic natural-language-inference
checkpoint: 330M parameters, one question-as-context template, premises
truncated to roughly the first 500 characters after tokenization, with 59
of the 191 pairs hitting the cap
(\path{paper-kbs/experiments/results/exp7_second_signal.json}). It scores
below the verdict (Hit@5
0.3125, MRR 0.1375), with inverted separation: gold median 0.0113 against
non-gold 0.05285 ($p = 0.120$). Truncation is not the explanation,
because six of the ten gold documents that fit the context window whole
still score below 0.01; nor is it an out-of-domain effect, as unrelated
topics still score corpus documents at up to 0.9973, so an NLI threshold
would abstain on nothing either. A span-level replay inverts harder: over
the 191 pairs and 3{,}489 span inferences at roughly 320 characters per
span (\path{paper-kbs/experiments/results/exp9_span_nli.json}), the max-pooled score is 0.4728 for gold against 0.91855 for non-gold
($p = 0.0074$), with Hit@5 0.1875. Rank fusion of the two signals repairs
neither (Hit@5 0.3125, MRR 0.111458). Replacing the judge checkpoint is
not the repair.

\textbf{What the failure tracks.}
Neither the generic NLI checkpoint, nor its span-level variant, nor rank
fusion over both reaches the verdict's own quality, which rules out the
checkpoint as the sole cause; the second corpus shows
the same gate losing nothing there (Section~\ref{sec:e4}), so the collapse
is associated with confuser density within one domain --- many
near-duplicate documents competing for the same question. That probe cannot
separate density from the other variables it moved: it changed the domain,
the language and the corpus size at the same time as it removed the
confuser clusters, so this paper reports an association that is confounded
by design rather than an elimination
\cite{zheng-llmjudge,panickssery-selfpref,lightman-verify,yin-knowwhat}.
The verdict earns its place as a filter and loses it as a ranker.

\begin{figure}[t]
\centering
\includegraphics[width=\textwidth]{fig5_e2e}
\caption{Agentic runs. (a)~Wall time per run; printed numbers are agent
turns. (b)~Citation gold-hit rate vs.\ mean agent turns (horizontal offsets added for
readability): the hybrid lane reaches fully grounded answers with the
fewest turns, the librarian lane spends the most turns for exhaustive
recall, and the bare agent cannot ground citations at all.}
\label{fig:e2e}
\end{figure}

\begin{figure}[t]
\centering
\includegraphics[width=0.5\textwidth]{fig4_inflation}
\caption{Evidence-criterion judge scores from the sixteen QDCVR-Hybrid tool
runs: all 191 judged candidates. Gold documents (green, $n = 17$; Q16's
gold never entered its candidate pool) and same-domain non-gold survivors
(orange, $n = 174$) overlap almost completely: medians 0.876 versus 0.863,
Mann--Whitney $p = 0.26$. The 0.5 acceptance threshold separates both
groups from zero, not from each other.}
\label{fig:inflation}
\end{figure}

\begin{figure}[t]
\centering
\includegraphics[width=\textwidth]{fig8_anatomy}
\caption{Anatomy of one real query (Q2; single logged QDCVR-Hybrid call).
Each line is one of the 12 fused candidates; the green band is the top-5
window the answer agent reads in full. (a)~The shipped re-ranking: the
judge scores five extruder-PLC telemetry reports $0.887$--$0.926$ and the
gold root-cause report $0.846$, demoting it from fusion rank~1 to
rank~8 --- outside the read window. (b)~The same call with the verdicts
used only as a filter: all 12 pass, the fusion order stands, and the gold
is the first document read. Same pool, same scores, same verdicts; only
the role of the induced ordering differs.}
\label{fig:anatomy}
\end{figure}

\subsection{E7: Non-functional behaviour}
\label{sec:e7}
\label{sec:robustness}\label{sec:eval-nonfunctional}

\textbf{Stability.}
A 27-run suite --- three task types $\times$ three lanes $\times$ three
rounds on the live multi-base instance
(\path{review/stability-multiround-20260930}) --- is summarised in
Table~\ref{tab:e7-stability}. Session-trace replay confirms 27/27
lane-contract compliance, with per-cell median wall clock spanning
34--76\,s. Caveats: correctness calls are the authors' trace readings,
cache state was uncontrolled, two platform cells re-ran once after a
timeout, and the $\pm$35\,s variance sits in the answer-generation
window.

\begin{table}[!htbp]
\centering
\small
\setlength{\tabcolsep}{2pt}
\caption{The 27-run stability suite (three task types $\times$ three lanes
$\times$ three rounds).}
\label{tab:e7-stability}
\begin{tabular}{p{2.4cm}p{6.4cm}p{4.3cm}}
\toprule
Task type & Expected behaviour & Outcome \\
\midrule
Answerable & correct attention formula & 9/9 correct \\
Not covered & report that the library holds no such content & 9/9 correct not-found \\
Out of corpus & refusal (the librarian lane stops at the catalog layer) & 9/9 refusals \\
\bottomrule
\end{tabular}
\end{table}

\begin{figure}[!htbp]
\centering
\includegraphics[width=\textwidth,height=0.80\textheight,keepaspectratio]{fig9_robustness}
\caption{The robustness tier at a glance. (a)~Paraphrase grid: per-method
Hit@5 on the original questions (blue) and their paraphrases (orange);
BM25's ceiling is vocabulary-dependent while dense, fusion, and the
filter repair hold. (b)~NLI separation over the 191 logged candidates
(span scores from 3{,}489 span inferences): both granularities
\emph{anti-separate} --- gold medians sit below non-gold medians.
(c)~Domain, language, and scale shift to 79 ingested arXiv papers: the
judged collapse
disappears when no same-family confuser clusters exist.
(d)~Graded abstention probes: the score the threshold must act on
stays high even where coverage is absent.}
\label{fig:robustness}
\end{figure}

\begin{table}[!htbp]
\centering
\scriptsize
\setlength{\tabcolsep}{2.5pt}
\caption{Robustness checks: the question each probe was designed to answer
and the measured result. All numbers come from released logs.}
\label{tab:robustness}
\begin{tabular}{p{2.7cm}p{3.9cm}p{5.9cm}}
\toprule
Check & Question & Result \\
\midrule
Cut-margin LOO & Is the 0.20 margin an artifact of selecting on all 16
questions? & The smallest-margin rule chosen on 15 questions transfers to
15/16 holdouts (miss = Q11, which the shipped 0.15 also loses); the
deployed 0.20 retains all 16 holdout golds. The $+0.02$ pad is
engineering judgment, not an LOO value. \\
Paraphrase grid ($n{=}16$) & Does the tool layer depend on the questions'
vocabulary? & BM25 1.00$\to$0.88 (MRR 0.958$\to$0.755); dense 0.94;
judged 0.50; filter = fusion 0.94. Pooled over 32 variants: 15 wins,
0 losses, 17 ties ($p < 0.0001$). Failure and repair are wording-robust. \\
Second signal (NLI) & Is inflation specific to the Laya checkpoint? &
A generic Chinese NLI ranks \emph{worse} (Hit@5 0.31, MRR 0.138;
$\rho = -0.13$) and \emph{anti-separates} (gold median 0.011 \emph{below}
non-gold 0.053, $p = 0.12$); out-of-corpus topics still score corpus
documents up to 0.997. Naive span-max NLI is worse still (0.19;
anti-separation $p = 0.007$); the rank-fusion ensemble recovers nothing
(0.31). \\
Symbolic check (Q6) & Can rule checks cover what model judges miss? &
The watch value ``31.35'' appears in the gold (1/1) and in no confuser
(0/11): the numeric case the judge mishandled is solved by a trivial
substring rule. \\
Graded probes ($n{=}6$) & Does the threshold behave differently between
partial and absent coverage? & No: the threshold keeps 10--12 of 12
candidates (up to 0.983) in every tier. Agent layer: scoped disclosures
on partial probes (one was reclassified as covered in a post-hoc human
read of that run --- the corpus held a
verified rework case); 3/3 absent probes returned explicit not-found. \\
Second corpus (100 arXiv papers) & Does the judged collapse survive a
domain, language, and scale shift with no confuser clusters? & No
collapse at 79 indexed papers (8{,}031 chunks): judged = fusion =
filter = 0.90 (one shared miss --- a gold paper that never became
retrievable).
The loss tracks confuser density, not domain, language, or size;
graph 0.20 without relation tagging (bulk import bypasses it). \\
Retrieval re-run (2026-10-12) & Do the six tool-layer methods still return
the gold with the model layer removed? & Yes, in 12 of 12 method--question
cells; modes 25.0/22.7, 25.0/10.2 and 37.3/33.7\,s (A, B, C); BM25 and RRF
0.0\,s, vector 0.1\,s
(\path{benchmark-suite/results/runs/run-20261012T062847Z-5f9faac/RETRIEVAL-COMPARE.md}).
This is the final pass of that run (the first pass returned no documents on
its vector arm). An $n = 2$ smoke check; no claim rests on it. \\
\bottomrule
\end{tabular}
\end{table}

Figure~\ref{fig:robustness} visualises the same checks: the cut-margin
variants, the anti-separating verdict and NLI scores, the shift to the
second corpus, and the graded abstention probes.

\textbf{Concurrency and integrity.}
A separate load campaign (\path{review/stress-20260923/STRESS-REPORT.md},
2026-09-23) wrote one library concurrently: 20 simultaneous document
creations finished in 0.8\,s, 20/20 HTTP 200, with no corruption of
\path{.tree-fs.json} or \path{.knowledge-base.yml} (list count 20);
20 simultaneous index builds finished in 4.1\,s (20/20, 620 chunks); 10
simultaneous experience creations finished 10/10. Lexical re-verification
found 3/3 probe documents and a vector top-3 check found one; the artifact
attributes that to homogeneous templates diluting the semantic signal, not
to data loss. This bounds concurrent writes only: no crash-mid-write or
partial-index adversarial test was run, so durability under interruption
remains untested.

\textbf{Integrity incidents.}
Two index-integrity incidents surfaced during evaluation setup --- one
document missing from its vector collection, one library whose collection
had been silently lost --- both caught by the pre-measurement checks and
repaired through the platform's own path (\texttt{index-document},
\texttt{reindex}).

\textbf{Deployment reality.}
The system is a working deployment, not a frozen benchmark, and its catalog
changes as campaigns run: the industrial corpus of
Section~\ref{sec:eval-setup} was measured while those five libraries were
resident, and the snapshot of 2026-10-12 shows the catalog afterwards, when
the two acquisition campaigns had populated five domain libraries holding
the 322 stored parts. Every figure here is therefore a timestamped
observation of a system in use, and none of them should be read as the
state of a static corpus.

\subsection{E8: Ablations of the verification gate and the search scope}
\label{sec:e8}
Every measurement above describes the system as shipped. Two matched
ablations test the two mechanisms that the design chapters argue for --- the
content gate of Section~\ref{sec:gate} and the content-routed partition of
Section~\ref{sec:org-placement} --- by disabling each in a control arm. Both
were run on the second
corpus --- ten questions over the five content-routed libraries, 322
stored parts --- with nothing else changed, and both return a largely null
result, which is what makes them informative about which layer actually
binds.

\textbf{Gate.}
Three arms shared one candidate pool (cap 12), one 4{,}000-character
evidence pack and one grader
(\path{paper-kbs/experiments/results/exp11_gate_ablation.json}): no gate,
the gate used as a filter, and a
diagnostic arm in which the answering channel received the judge's own
segment-level evidence pack. Gold was packed at rank 1 in all eight
gold-answerable questions in all three arms, so the gate had no ordering
work to do: it removed no gold document (8/8 packed gold survived), and on
nine of ten questions the gated arms received exactly the documents the
ungated arm received. The only cell the gate shortened was QK09 (12 to 4
documents), which changed no grade. The gated arms therefore differ from the
control only where the answering model's own output varied: QK01
($0\to5$) and QK03 ($3\to0$), giving mean grades over the
gold-answerable questions of 1.00 and 1.25 (1.80 and 2.00 over all ten;
Table~\ref{tab:ablations}a). Two cells moving in opposite directions carry
no statistical weight (sign test $p = 1.0$, $n = 10$), and QK01 is the
clearest evidence that the swing is answering-model variance: both of its
arms received the same twelve documents and the same 4{,}000-character
budget and graded 0 and 5. The segment-granularity arm scored 0.375 on the
same subset, below the control, losing three answers the document-level gate
kept and gaining one. What the ablation
does show is where the remaining failures are. Four of the eight
gold-answerable questions (QK02, QK04, QK05, QK08) scored zero in every
arm and a fifth scored partial credit in one arm and zero in the other
two, even though the gold document was packed at rank 1 in all of them.
Both non-gold questions --- an unanswerable and an out-of-corpus
probe --- were answered correctly as abstentions in every arm. The binding
constraint on this corpus is the reading step, not candidate ordering,
which is the same layer at which the industrial corpus locates honest
abstention (Section~\ref{sec:e5}).

\textbf{Search scope.}
The second ablation varies what the retriever may see rather than what
survives it: the library that holds the gold, a different library, all
five libraries at once, and the catalog lane alone
(Table~\ref{tab:ablations}b;
\path{paper-kbs/experiments/results/exp12_partition_ablation.json}). Pinned to the gold library, Hit@5 is 1.0
with the gold at rank 1 among 2.8 candidates. Pinned to a different
library, Hit@5 is 0.0 --- and the wrong library still yields 6.9
surviving candidates, so the failure is confident and silent rather than
an empty result. Searching all five libraries at once reproduces the
gold-pinned result exactly (1.0, gold at rank 1, 2.5 candidates), so
partitioning costs no recall on this corpus; the catalog path also
reaches 1.0 but at a mean gold rank of 1.75 and only after the gate has
reduced 236 candidates to 11.4. Catalog routing selected the correct
library first for 6 of the 8 questions. A misplacement is therefore not a
ranking penalty but a potential total recall loss --- and because every
stored filename carries the acquisition prefix
(\texttt{\{domain\}\_\_\{arxiv\_id\}\_\_\{slug\}}, present on 322/322
parts), name-based routing is not an independent measurement: what the
platform demonstrates is routing coverage, not routing accuracy. Both
ablations are single-corpus, single-engine, ten-question measurements and
are reported as bounds rather than as effect sizes. One protocol deviation
is material: the gate ablation was rebuilt around the judging tool applied
to a frozen shared candidate set, because the configuration switch named in
earlier platform documentation does not exist in the shipped tool layer.
The ungated arm is therefore a fusion-ordered control rather than a replay
of a formerly shipped setting, and only the within-run contrast is
comparable. Two of the thirty gate cells were re-run after an
authentication failure terminated them mid-run; the replacements are
recorded alongside the originals.

\begin{table}[t]
\centering
\small
\caption{Component ablations on the second corpus (ten questions, 322
stored parts). (a) Verification gate: identical candidate pool, evidence
pack and grader in all three arms; the third arm delivers the judge's own
segment-level evidence to the answering channel and is a diagnostic rather
than a suppression policy. Grades are means over the eight
gold-answerable questions and over all ten. (b) Search scope: the same
queries under four retrieval scopes; ranks and counts are means over the
eight gold-answerable questions.}
\label{tab:ablations}
\begin{tabular}{@{}>{\raggedright\arraybackslash}p{3.5cm}rrrr@{}}
\toprule
\multicolumn{5}{@{}l}{\textbf{(a) Verification gate}}\\
\midrule
Arm & Grade (gold) & Grade (all) & Docs packed & Latency (s)\\
\midrule
No gate & 1.00 & 1.80 & 12.0 & 6.9 \\
Gate as filter & 1.25 & 2.00 & 11.2 & 7.7 \\
Gate at segment granularity & 0.375 & 1.30 & 11.2 & 6.4 \\
\bottomrule
\end{tabular}

\vspace{3pt}
\begin{tabular}{@{}lrrrr@{}}
\toprule
\multicolumn{5}{@{}l}{\textbf{(b) Search scope}}\\
\midrule
Scope & Hit@5 & Gold rank & Candidates & Kept after gate\\
\midrule
Library holding the gold & 1.0 & 1.00 & 2.8 & 2.8 \\
A different library & 0.0 & --- & 7.1 & 6.9 \\
All five libraries & 1.0 & 1.00 & 2.5 & 2.5 \\
Catalog lane only & 1.0 & 1.75 & 235.8 & 11.4 \\
\bottomrule
\end{tabular}
\end{table}

\section{Discussion and lessons learned}
\label{sec:discussion}

The evidence in Section~\ref{sec:eval} supports a distinction that the
retrieval literature tends to blur: finding a document and being allowed to
cite it are different operations, and the second one is a system property. The
lessons below are ordered from the retrieval layer outwards, because that is
the order in which we learned them.

\noindent\textbf{L1. Vectors are fast; content is accurate, and verification
has two costs.} The tool-layer comparison is the cleanest statement of the
retrieval subproblem: on a curated corpus a millisecond-level BM25 beats an
18\,ms dense retriever, and neither knows whether a document \emph{answers}
the question. The gate converts ranked candidates into accountable evidence
with per-citation provenance. Its costs are temporal (about 0.5\,s per dozen
candidates end to end; 1--2\,s per six-document batch for the resident worker)
and, on this corpus, accuracy-related: 0.94\,$\to$\,0.44 Hit@5 under
same-domain inflation when the verdicts are used as an ordering. Auditability
is what the gate buys today; ranking quality is what calibration must earn
back.

\noindent\textbf{L2. Latency engineering is turn engineering.} With retrieval
sub-second, the only lever that matters for interactive latency is the number
of agent turns. Tool whitelists that remove detour tools, one-call server-side
pipelines, and small judge batches that avoid transport timeouts together took
the hybrid lane from about 29 turns in the earliest serial-skill configuration
to 3--5 turns in the present measurements. The same accounting explains the
engine comparison: the third engine's wall clock exceeded its engine-side time
by 50--90\,s per call, which is harness startup, not retrieval.

\noindent\textbf{L3. The verifier is a model too, so calibrate it.} The
inflation analysis is the paper's cautionary result: a verifier inherits
domain-internal blind spots, and a threshold calibrated on cross-domain
negatives (0.5) does not separate same-domain near-misses. Fail-closed
semantics and forced reading contain the damage; they do not replace
calibration. The second-signal probes sharpen the point: a generic entailment
model fails too, but differently, anti-separating golds while still rejecting
no out-of-corpus topic, and a symbolic rule decides the numeric case that the
shipped judge mishandles. Verifier choice changes which errors you get, not
whether you get any.

\noindent\textbf{L4. Honesty needs an enforcement layer, not a prompt.} Lane
separation works because the server physically withholds tools: a
librarian-lane agent cannot call the vector tool even if instructed to. The
out-of-corpus probes show the same asymmetry for abstention. A threshold
cannot refuse; it accepted every out-of-corpus candidate it was shown. The
reading-and-answer layer refused all six probe prompts with disclosed searches
and zero fabricated citations. A bare agent without the platform refused
explicitly in two of four runs, and in the others produced confident,
unsourced, partly wrong answers whose wrongness no audit trail exposes.
Unassisted honesty is a coin flip; enforced honesty is inspectable.

\noindent\textbf{L5. Production incidents taught honest observability.} Three
early incidents shaped the implementation: a corrupted index segment that made
one library silently un-retrievable across lanes, an out-of-memory indexing
run that returned success with an empty payload, and judge failures that had
to reject candidates rather than pass them. The common thread is that every
layer reports what it actually verified and every consumer treats the
unverifiable as failed. Table~\ref{tab:failures} consolidates the incident log
into the failure-mode catalogue we wish we had on day one, extended with the
three defects the organization and functional campaigns surfaced.

\begin{table}[!htbp]
\centering
\scriptsize
\renewcommand{\arraystretch}{0.95}
\setlength{\tabcolsep}{3pt}
\caption{Failure-mode catalogue from production and campaign use: how each
failure presents, the detection signal, the fix, and the residual risk an
adopter still carries.}
\label{tab:failures}
\begin{tabular}{p{3.2cm}p{3.2cm}p{3.2cm}p{3.2cm}}
\toprule
Failure & Presents as & Detection signal & Fix / residual risk \\
\midrule
Corrupted index segment (silent loss) & ``Success, zero results'' for one
library across lanes & Treat zero-result success as suspect; cross-lane
census & Re-index the collection; residual: any new silent-loss mode needs
the same tripwire \\
Out-of-memory during indexing & Success envelopes with empty vector payloads
& Judge success by chunk counts, not envelopes & Memory hygiene between
embedding and judging workloads; residual: other silent truncations \\
Driver contention (Windows) & Embedding latency 68\,s $\to$ 27\,min &
Latency budget on the embedding path & Isolate embedding and judging
workloads; on Linux, watch VRAM contention instead \\
Judge worker cold start & First judge batch stalls 10--12\,s & Worker
residency heartbeat & Keep the resident worker warm; residual: first-call
timeouts still fail closed \\
Un-whitelisted harness tools & 14--170\,s detours per run & Turn count and
tool-call audit log & Server-side per-lane whitelists; residual: new harnesses
need new whitelists \\
Judge or engine errors & Wrong or missing verdicts & Fail-closed rejects at
the gate & Errors reject, never pass; residual: availability, not
correctness \\
Rename without re-index & Chunk count 18 where 9 exist; the document matches
its own former path & Three-source audit; near-duplicate report &
Force re-index; residual: the write path still cannot re-index implicitly \\
Graph node not created & A library reports two of three documents &
Per-library document count versus graph overview & Rebuild the graph for that
document; residual: construction is not yet guaranteed by ingestion \\
Task record outlives its task & A finished build still reports ``running'' &
Compare task status against the artifact it was to produce & Query the
artifact, not the task; residual: status is not completion evidence \\
\bottomrule
\end{tabular}
\end{table}

\noindent\textbf{L6. The acquisition gate sets the ceiling for everything
above it.} In the 100-document campaign, 21 submitted documents never became
retrievable --- the campaign record does not separate parse losses from
record-keeping losses, and the material is dense scanned work with formulas
and tables. The system's value here is
not that the rate is low, but that a loss becomes a reported outcome: a
document that fails a gate and is nevertheless stored as empty or garbled
would enter the
catalog, receive a description, be indexed, and be retrieved and cited by the
same machinery that cites good documents, with no signal anywhere that
anything was lost. The same argument applies to a missing index or a stale
description. This is why the acquisition layer owns a post-condition rather
than a warning, and why the knowledge layer is a component of the system
rather than a precondition to retrieval.

\noindent\textbf{L7. Organization is a retrieval budget.} Content routing
decides what the candidate generator can see. A catalog that is curated and
consistent lets an agent choose a library and a document set from metadata
alone, which is the mechanism behind the catalog-first lane and behind
cross-library discovery in a deployment with nineteen libraries. An
ungoverned catalog produces the opposite: candidate generation degenerates
into a global vector search whose recall is limited by the cap, and no
downstream verification can recover a document that was never proposed. The
organization layer is therefore not administrative overhead around retrieval;
it is where retrieval's recall is decided.

The lessons share one implication for how systems of this kind should be
evaluated. A claim about enforcement is a claim about a control, so it needs a
matched ablation: the same pipeline with the mechanism disabled, the same
questions, the same corpus. Where we have that control (retrieval with and
without verification at the tool layer) the claim is quantified; where we did
not have it at first (an end-to-end agent arm with the gate disabled) we fixed
the protocol before running it rather than after
(Section~\ref{sec:eval-agent}). The verifier's negative result is the reason
we insist on the distinction: without the matched ablation, a filter and a
re-ranker are indistinguishable from the outside, and only one of them is
defensible.

\section{Threats to validity}
\label{sec:threats}

\noindent\textbf{Scale.} The retrieval corpus is 42 documents across five
libraries; the campaigns scale it to 322, and a second corpus adds 79. Absolute
Hit@5 values (BM25 at 1.00 on the industrial corpus) do not transfer to much
larger collections, where the hybrid candidate cap and the librarian's shelf
pruning become load-bearing. The shipped benchmark suite is the scaling path; a
hundreds-of-documents run with an author-independent question set remains the
full answer.

\noindent\textbf{Question set.} Sixteen tool-layer and four agent-layer
questions on the industrial corpus, ten on the second corpus and forty cloze per
language in the campaign were developer-authored. Gold documents and fact
markers are released, but the samples are small and favour the corpora's
vocabulary. Q16's question embeds its gold document's filename token,
guaranteeing lexical success; the fix is scheduled, and because Q16 is the one
question on which the fusion stage loses a gold, the correction can only tighten
the reported decomposition. The arXiv documents are named
\texttt{\{domain\}\_\_\{identifier\}\_\_\{slug\}}, so routing accuracy against
that prefix is not an independent measurement, and we report coverage rather
than routing accuracy (Section~\ref{sec:e2}).

\noindent\textbf{Experimental controls.} The tool layer isolates the gate:
among its nine configurations are retrieval without verification and the same
retrieval with the gate in either role, so the ranking-layer claim rests on a
matched control. The agent layer lacked one: its four arms are three grounded
lanes (retrieval \emph{and} verification) plus a bare agent (neither), which
cannot separate them. The missing arm, a grounded lane with the gate disabled
and the fusion order returned, was added late in revision, its protocol fixed
before the run (Section~\ref{sec:e8}), plus a second ablation varying
the search partition rather than the gate (\S\ref{sec:e8}). Both inherit
the same questions, libraries, engine and turn budget; neither was tuned after
seeing its results.

\noindent\textbf{Evidence provenance is uneven.} The acquisition, organization
and functional campaigns are audited runs with released per-check outcomes. The
experience subsystem rests on case evidence: a distillation campaign with
thirteen drafts, one sedimented experience, and an approval path never
exercised. The 100-document ingestion is a derived state file, not a timed
per-document log, and its accounting is not closed: 100 documents were
submitted, 91 names appear in the derived state, and 79 are present in the
library. The whole gap is unattributed: no artifact records a parse status
per paper, so the 21 missing papers may have failed parsing, may have failed
a later gate, or may never have been recorded. A 21\% retrievability gap is
therefore a bounded estimate of the parse-failure rate, not a measurement of
it, and no claim in this paper depends on the distinction. A system paper's
most useful contribution may be the honesty
of its ledger.

\noindent\textbf{Judge and engine dependence.} Scores and turn counts use one
judging checkpoint and one agent engine per campaign; lanes rather than models
are compared, but absolute numbers move with engine versions. The deployed gate
is still a single checkpoint, although the two strongest verifier findings are
no longer single-signal: a generic entailment model also fails (it
anti-separates golds and rejects no out-of-corpus topic), and a symbolic rule
decides the numeric case the shipped judge mishandles. Trained
evidence-verification models and learned signal ensembles remain open; the naive
span-max probe and the rank-fusion ensemble answer the two easiest upgrades
negatively.

\noindent\textbf{Design surface.} Three lanes triple the policy surface an
adopter must test. The evidence supports two operating points: the catalog-first
lane as the exhaustive option, the orchestrated lane as the latency-tuned
default. A two-lane deployment is defensible today; query-shape auto-selection
is future work.

\noindent\textbf{Consumer-side evidence.} The honesty contract is evaluated on
the system only, that not-found reports are well formed and non-fabricated.
Whether an operator acts on a not-found report, escalates to a human, or ingests
a recommended source is not measured: the paper's main missing external-validity
leg. A study is a small operator trial over a fixed probe set with two
conditions, platform reports versus the same model's unaided answer, scored on
the downstream action taken; probe set and transcripts are released so it can be
run independently.

\noindent\textbf{Nondeterminism.} Agent runs are stochastic; we report two
repetitions of the sixteen-cell grid (citation agreement 16/16, with wall time
and turn counts fluctuating substantially and not a distribution), retry only
zero-turn transport flakes, and release all per-run logs including timeout
re-runs. Fact-marker checks are automated substring matches; we additionally
report decontaminated checks computed after removing answer lines that name a
gold document, and the out-of-corpus answers were audited manually.
Deterministic layers are bit-identical across two rounds except for latency
fields.

\noindent\textbf{Untested surfaces.} Three designed surfaces carry no
controlled evidence and no claim is made about them. Library merge and split
are described as life-cycle transitions but were never executed under
controlled conditions, and neither is a single tool call: both are operator
procedures over the document, catalog and index tools
(Section~\ref{sec:e2}, Table~\ref{tab:org-ops}). The experience store's decay
and scheduled-distillation paths were not exercised end to end
(Section~\ref{sec:e3}). Durability under a crash mid-write was not tested: the
write ordering and the visible index post-condition are as designed, and the
concurrency campaign covers interleaved writes, not process death
(Section~\ref{sec:e7}). Authentication and token management exist in the
deployment and are outside this evaluation; the corpora used here are
author-owned and were processed on a single workstation, so no privacy or
adversarial-access claim is made either.

\section{Conclusion}
\label{sec:conclusion}

QDCVR turns ``the retriever found it'' into ``the system verified it''. Its
design position is that the governed object of a knowledge-base management
system is content: content decides where a document belongs and whether a
candidate may be cited, and the absence of content that answers a question is
itself an outcome the system must report. That position is implemented as
contracts rather than conventions (RQ1): parse gates that reject visibly,
placement rules that read the document, a three-source consistency audit over
disk, tree registry and library metadata, tool whitelists enforced server-side,
and verdicts that fail closed.

The five research questions are answered by measurement rather than
assertion. Acquisition and content routing hold at the scale we tested, while
21 of the 100 submitted papers never became retrievable and the gap between the
submitted, recorded and retrievable counts is reported as an observability gap
rather than explained away (RQ2). Staged, multi-channel retrieval with content
adjudication reaches documents that one-shot vector top-$k$ misses, and the
benefit comes from staged recall plus content verification, not from the fusion
of the two retrieval stages (RQ3). Verification costs well under a second and
buys per-citation provenance, but as an ordering it inherits a same-domain
blind spot that we quantify and partially repair; the failure tracks confuser
density rather than the checkpoint, domain, language or scale (RQ4). Honest
abstention lives in the reading and answering layer, never in the threshold;
end-to-end latency belongs to model turns; and engine choice moves cost by two
orders of magnitude at the same quality tier (RQ5).

The roadmap follows from the same ledger. Verifier calibration against
per-library null distributions, trained evidence-verification models, and the
verdict-as-filter role that the tool layer already supports are the immediate
next steps; larger-corpus evaluation with an author-independent question set,
consumer-side evidence that not-found reports serve their readers, and lane
auto-selection from query shape are the open empirical ones. The direction is
towards retrieval systems whose trustworthiness is an inspectable property of
the pipeline rather than a hope about the model.

\section*{Data availability}
The platform source code, the acquisition and organization campaign logs, the
tool-layer and agentic experiment scripts, the gold question set with fact
markers, the paraphrase grid, the leave-one-out margin replay, the
second-signal and span-level NLI replays, the graded abstention probes, the
cross-lingual acquisition results, the gate and partition ablation runs, and
all raw per-run JSON logs (tool-layer retrieval rows and agentic transcripts,
including the timeout re-runs) are available at
\url{https://github.com/kingdol666/rag-knowledge} (benchmark suite:
\texttt{benchmark-suite/}; paper artifacts and evidence inventories:
\texttt{paper-kbs/}).

\section*{Ethics statement}
This article does not report studies involving human participants, human
data, or animals.

\section*{Declaration of generative AI and AI-assisted technologies in the manuscript preparation process}
During the preparation of this work the authors used LLM-based coding
agents to assist with experiment execution, data tabulation, figure
preparation, and manuscript drafting. After using this tool/service, the
authors reviewed and edited the content as needed and take full
responsibility for the content of the published article. All experimental
numbers originate from logged runs of the deployed system on the authors'
machine.

\section*{Declaration of competing interest}
The authors declare that they have no known competing financial interests
or personal relationships that could have appeared to influence the work
reported in this paper.

\section*{CRediT authorship contribution statement}
\textbf{Author 1:} Conceptualization, Methodology, Software, Investigation,
Data curation, Writing -- original draft. \textbf{Author 2:} Validation,
Writing -- review \& editing, Supervision. %% AUTHOR ACTION: adjust.

\section*{Funding}
This research did not receive any specific grant from funding agencies in
the public, commercial, or not-for-profit sectors. %% AUTHOR ACTION

\section*{Acknowledgments}
The authors thank the maintainers of the open-source components the platform
builds on. %% AUTHOR ACTION: adjust or remove.

\appendix
\setcounter{table}{0}
\section{Per-question tool-layer results}
\label{app:perq}
Table~\ref{tab:appendixranks} breaks the tool-layer comparison down to
individual questions: the best gold rank each configuration reached in
its top-5 window, for every gold-annotated question. Reading the columns
side by side localizes every aggregate number in the paper: BM25 and RRF
miss at most one question (Q16's confuser family and Q7's telemetry
cluster account for the remaining sub-1.0 MRRs); the judged hybrid's nine
misses are eight questions where same-domain confusers outranked
the gold (Q2, Q4, Q5, Q6, Q8, Q9, Q11, Q12) plus Q16, where the confuser
family crowded the gold out upstream;
and judge-as-filter restores the gold on every one of the eight
judging losses while inheriting Q16's upstream failure --- the
decomposition that Fig.~\ref{fig:marginvariants}(b) aggregates and
Fig.~\ref{fig:anatomy} dissects one query at a time.

% Auto-generated by experiments/derive_appendix.py -- do not edit.
\begin{table}[t]
\centering
\scriptsize
\setlength{\tabcolsep}{5.5pt}
\caption{Per-question best gold rank (top-5 window; `--' = gold absent) for the seven tool-layer pipelines (logged runs) and the two derived configurations (offline replay of the same logged QDCVR-Hybrid pools as Fig.~\ref{fig:marginvariants}). Fusion = fusion-only (no judge); Hybrid = the shipped QDCVR-Hybrid judged ordering; Filter = judge-as-filter. Question numbers and texts: Table~\ref{tab:questions}.}
\label{tab:appendixranks}
\begin{tabular}{lccccccccc}
\toprule
Q & BM25 & Dense & RRF & Graph & 2-stage & Fusion & Hybrid & Filter & RRF-fuse \\
\midrule
1 & 1 & 1 & 1 & 1 & 1 & 1 & 4 & 1 & 1 \\
2 & 1 & 1 & 1 & -- & 1 & 1 & -- & 1 & 1 \\
3 & 1 & 1 & 1 & -- & 1 & 1 & 1 & 1 & 1 \\
4 & 1 & 1 & 1 & -- & 1 & 1 & -- & 1 & 1 \\
5 & 1 & 1 & 1 & 1 & 1 & 1 & -- & 1 & 2 \\
6 & 1 & 1 & 1 & 1 & 1 & 1 & -- & 1 & -- \\
7 & 3 & 5 & 3 & -- & 3 & 5 & 3 & 5 & 3 \\
8 & 1 & 1 & 1 & 1 & 1 & 1 & -- & 1 & 2 \\
9 & 1 & 1 & 1 & 1 & 1 & 1 & -- & 1 & 3 \\
10 & 1 & 1 & 1 & 1 & 1 & 1 & 3 & 1 & 1 \\
11 & 1 & 2 & 1 & 1 & 1 & 2 & -- & 2 & 5 \\
12 & 1 & 1 & 1 & 1 & 1 & 1 & -- & 1 & 3 \\
13 & 1 & 1 & 1 & -- & 1 & 1 & 2 & 1 & 1 \\
14 & 1 & 1 & 1 & 1 & 1 & 1 & 1 & 1 & 1 \\
15 & 1 & 1 & 1 & -- & 1 & 1 & 5 & 1 & 2 \\
16 & 1 & -- & 2 & -- & 2 & -- & -- & -- & -- \\
\bottomrule
\end{tabular}
\end{table}


\section{Supplementary tables}
\label{app:tables}
The tables below collect reference material that supports the design and
evidence chapters but is not required to follow the argument: the storage
data model, the cross-lingual acquisition result, the gold question set and
the lifecycle-operation evidence.
\begin{table}[!htbp]
\centering
\footnotesize
\setlength{\tabcolsep}{3.5pt}
\caption{Stores and their principal contents. Field names are the ones used
by the implementation and by the audit tools.}
\label{tab:datamodel}
\begin{tabular}{>{\raggedright\arraybackslash}p{3.4cm}>{\raggedright\arraybackslash}p{9.7cm}}
\toprule
Store & Contents \\
\midrule
Tree registry (\path{.tree-fs.json}) & version, root identifier and name,
timestamps, node counter and aggregate statistics; per node:
identifier, name, type, parent, path, children, metadata map, timestamps.
Authoritative structure of the collection. \\
Library metadata (\path{.knowledge-base.yml}) & library identifier, name,
description, creation and update timestamps, per-library configuration
(including distillation settings); per document: path, name, description,
tags, content hash, file size, vector-index metadata, free-form metadata.
The retrieval index read by catalog search. \\
Document files & Markdown body per document or per split part, plus an image
directory for figures extracted during parsing. \\
Vector store & one collection per library (cosine space, 1024-dimensional
embeddings, HNSW index), holding segment embeddings of 500 characters with
50-character overlap. \\
Graph store & nodes for documents, libraries, tags, folders and entities;
edges for membership, tagging, relatedness, identity and containment. \\
Task and session databases & pollable task records for parsing, indexing,
graph builds and distillation; chat sessions; service credentials;
distillation signals. \\
\bottomrule
\end{tabular}
\end{table}

\begin{table}[t]
\centering
\caption{Cross-lingual ingestion, 400 Wikipedia pages per language,
first round. The second round reproduced these values bit-for-bit, so
one column is shown. Self-retrieval Hit@1 uses 40 queries per
language.}
\label{tab:acq-crosslang}
\small
\setlength{\tabcolsep}{3pt}
\begin{tabular}{@{}lcccc@{}}
\toprule
\textbf{Language} & \textbf{Ingestion} & \textbf{KB membership} & \textbf{Storage} & \textbf{Self-retrieval} \\
 & \textbf{success} & \textbf{accuracy} & \textbf{completeness} & \textbf{Hit@1} \\
\midrule
English (en) & 1.0 & 0.875  & 0.975   & 0.975 \\
Chinese (zh) & 1.0 & 1.0    & 0.9559  & 1.0 \\
Japanese (ja) & 1.0 & 1.0    & 0.9402  & 1.0 \\
\bottomrule
\end{tabular}
\end{table}

\begin{table}[!htbp]
\centering
\scriptsize
\setlength{\tabcolsep}{3pt}
\caption{Gold-annotated question set over the five production knowledge
bases. ``Gold'' lists unique filename fragments of the ground-truth
documents. All sixteen questions serve the tool-layer experiment of
Section~\ref{sec:e4}; Q2, Q4, Q6, Q7 also
serve the agentic experiment of Section~\ref{sec:e5}.}
\label{tab:questions}
\begin{tabular}{lp{5.3cm}p{4.0cm}p{2.8cm}}
\toprule
\# & Question (abridged, translated) & Gold document(s) & Design intent \\
\midrule
Q1 & EVA encapsulation film lamination: how to set the temperature, pressure
and time window; how to control cure degree & EVA process essentials; EVA
hydro-thermal aging mechanism & multi-gold \\
Q2 & EL inspection shows edge delamination: root cause and corrective
actions & EL-2025-0317 delamination RCA report & incident $\to$ RCA \\
Q3 & Why does humid aging of EVA generate acetic acid and corrode ribbons;
moisture transport mechanism & EVA hydro-thermal aging mechanism & mechanism \\
Q4 & Moisture-barrier behavior of POE encapsulant; lamination differences
vs.\ EVA & POE process essentials & cross-material \\
Q5 & Summer power drop at a PV plant: systematic troubleshooting order &
Power-drop troubleshooting manual; O\&M lessons & synonym gap
(``drop'' vs.\ ``plunge'') \\
Q6 & Trigger condition, nominal value and tolerance of the line-1027 part
weight alarm & Watch-clause document (31.35\,g $\pm$0.05\,g, 5\,min) &
numeric exactness \\
Q7 & Melt temperature fell 206$\to$180$^\circ$C on extruder line 1: final
root-cause verdict & TC0 dead-loop RCA revision (among 20 telemetry dups) &
confuser density \\
Q8 & PET biaxial stretching process flow and critical temperature ranges &
PET process regulation & plain factual \\
Q9 & Frequent inverter alarm trips: handling steps & Inverter alarm handling
manual & PV O\&M family \\
Q10 & Hot-spot acceptance criteria and disposal & Hot-spot detection and
disposal procedure & PV O\&M family \\
Q11 & String mismatch causes; IV-curve diagnosis & String mismatch analysis
method & PV O\&M family \\
Q12 & PP cast-film extrusion parameters; chill-roll temperature & PP cast
film process regulation & plain factual \\
Q13 & Weight closed-loop run1: staircase excitation design & Weight
optimization staircase record & telemetry cluster \\
Q14 & FT full-function regression after P0 optimization & FT regression test
record & telemetry cluster \\
Q15 & Multimodal diagnosis: fusing scalar, vector and image signals &
Three-modality closed-loop record & telemetry cluster \\
Q16 & Post-compensation recovery verification of the heatdecay092 event &
heatdecay092 compensation verification (among 3 near-dup incident docs) &
confuser density \\
\bottomrule
\end{tabular}
\end{table}

\begin{table}[t]
\centering
\small
\setlength{\tabcolsep}{4pt}
\caption{Lifecycle operations exercised on disposable test libraries.
Merge and split were analysed but never executed.}
\label{tab:org-ops}
\begin{tabular}{@{}l>{\raggedright\arraybackslash}p{11.2cm}@{}}
\toprule
Operation & Observed outcome \\
\midrule
Create & Pass, document indexed and retrievable \\
Rename & Pass, stale chunks remained until forced re-index (18 recorded, 9 expected) \\
Move & Pass, source 1$\rightarrow$0$\rightarrow$1, target 0$\rightarrow$1$\rightarrow$0 \\
Delete & Pass, batch delete 1/1; both test libraries removed \\
Re-index & Automatic and batch paths did not purge; pass with \texttt{force=true} (18$\rightarrow$9) \\
Merge, split & Not evaluated \\
Task status & Fail, reports \texttt{running} after \texttt{finished\_at} is set \\
Update check & Fail, \texttt{ragctl version} timed out after 30\,s \\
\bottomrule
\end{tabular}
\end{table}



\bibliographystyle{unsrtnat}
\bibliography{refs,refs_new_verified}

\end{document}
